\documentclass[12pt,a4paper]{article}

% ============== Основные пакеты ==============
\usepackage[utf8]{inputenc}
\usepackage[T2A]{fontenc}
\usepackage[russian]{babel}

\usepackage{graphicx}    % для логотипа и рисунков
\usepackage{geometry}    % поля
\usepackage{setspace}    % межстрочный интервал
\usepackage{indentfirst} % отступ первой строки
\usepackage{fancyhdr}    % колонтитулы
\usepackage{caption}     % подписи к рисункам

\usepackage[hidelinks]{hyperref}
\usepackage{amsmath}
\usepackage{float}
\usepackage{amssymb}
\usepackage{enumitem}

\usepackage{xcolor}
\usepackage{listings}


% ============== Параметры страницы ==============
\geometry{
  left=30mm,
  right=15mm,
  top=20mm,
  bottom=20mm,
}

\onehalfspacing         % полуторный интервал
\parindent=1.25cm       % абзацный отступ

% ============== Колонтитулы и номера страниц ==============
\pagestyle{fancy}
\renewcommand{\headrulewidth}{0pt}
\fancyhf{}
\fancyfoot[C]{\thepage} % номер страницы по центру внизу

% ============== Оформление подписей рисунков ==============
\usepackage{caption}
\captionsetup[figure]{%
  labelfont=bf,
  % убрали labelsep=emdash
  labelsep=space,         % обычный пробел
  justification=centering,
  font=small,
  skip=0.5\baselineskip,
}

\definecolor{codebg}{RGB}{253,253,253}
\definecolor{keywordcolor}{RGB}{0,0,205}
\definecolor{stringcolor}{RGB}{34,139,34}
\definecolor{commentcolor}{RGB}{120,120,120}
\definecolor{functioncolor}{RGB}{199,21,133}
\definecolor{numbercolor}{RGB}{255,140,0}
\definecolor{builtincolor}{RGB}{0,102,204}
\definecolor{typecolor}{RGB}{128,0,128}

\lstdefinelanguage{PythonVSCode}{
    language=Python,
    morekeywords={
        import, from, as, class, def, return, if, elif, else,
        while, for, in, try, except, with, lambda, yield, True, False, None
    },
    morekeywords=[2]{np, plt, pd, sp, sympy, torch},
    morekeywords=[3]{linspace, arange, array, tensor, mean, std, sum,
                     plot, show, figure, xlabel, ylabel, title, legend, grid},
    sensitive=true,
    morecomment=[l]{\#},
    morestring=[b]',
    morestring=[b]"
}

% Стиль lstlisting с русской поддержкой
\lstset{
    language=PythonVSCode,
    backgroundcolor=\color{codebg},
    basicstyle=\ttfamily\small,
    keywordstyle=\color{keywordcolor}\bfseries,
    keywordstyle={[2]\color{builtincolor}},
    keywordstyle={[3]\color{functioncolor}},
    stringstyle=\color{stringcolor},
    commentstyle=\itshape\color{commentcolor},
    numberstyle=\tiny\color{gray},
    numbers=left,
    numbersep=10pt,
    showstringspaces=false,
    breaklines=true,
    frame=none,
    tabsize=4,
    captionpos=b,
    escapeinside={(*@}{@*)},
    inputencoding=utf8,
    extendedchars=true,
    literate=
        {а}{{\selectfont\char224}}1 {б}{{\selectfont\char225}}1 {в}{{\selectfont\char226}}1
        {г}{{\selectfont\char227}}1 {д}{{\selectfont\char228}}1 {е}{{\selectfont\char229}}1
        {ё}{{\selectfont\char184}}1 {ж}{{\selectfont\char230}}1 {з}{{\selectfont\char231}}1
        {и}{{\selectfont\char232}}1 {й}{{\selectfont\char233}}1 {к}{{\selectfont\char234}}1
        {л}{{\selectfont\char235}}1 {м}{{\selectfont\char236}}1 {н}{{\selectfont\char237}}1
        {о}{{\selectfont\char238}}1 {п}{{\selectfont\char239}}1 {р}{{\selectfont\char240}}1
        {с}{{\selectfont\char241}}1 {т}{{\selectfont\char242}}1 {у}{{\selectfont\char243}}1
        {ф}{{\selectfont\char244}}1 {х}{{\selectfont\char245}}1 {ц}{{\selectfont\char246}}1
        {ч}{{\selectfont\char247}}1 {ш}{{\selectfont\char248}}1 {щ}{{\selectfont\char249}}1
        {ъ}{{\selectfont\char250}}1 {ы}{{\selectfont\char251}}1 {ь}{{\selectfont\char252}}1
        {э}{{\selectfont\char253}}1 {ю}{{\selectfont\char254}}1 {я}{{\selectfont\char255}}1
        {А}{{\selectfont\char192}}1 {Б}{{\selectfont\char193}}1 {В}{{\selectfont\char194}}1
        {Г}{{\selectfont\char195}}1 {Д}{{\selectfont\char196}}1 {Е}{{\selectfont\char197}}1
        {Ё}{{\selectfont\char168}}1 {Ж}{{\selectfont\char198}}1 {З}{{\selectfont\char199}}1
        {И}{{\selectfont\char200}}1 {Й}{{\selectfont\char201}}1 {К}{{\selectfont\char202}}1
        {Л}{{\selectfont\char203}}1 {М}{{\selectfont\char204}}1 {Н}{{\selectfont\char205}}1
        {О}{{\selectfont\char206}}1 {П}{{\selectfont\char207}}1 {Р}{{\selectfont\char208}}1
        {С}{{\selectfont\char209}}1 {Т}{{\selectfont\char210}}1 {У}{{\selectfont\char211}}1
        {Ф}{{\selectfont\char212}}1 {Х}{{\selectfont\char213}}1 {Ц}{{\selectfont\char214}}1
        {Ч}{{\selectfont\char215}}1 {Ш}{{\selectfont\char216}}1 {Щ}{{\selectfont\char217}}1
        {Ъ}{{\selectfont\char218}}1 {Ы}{{\selectfont\char219}}1 {Ь}{{\selectfont\char220}}1
        {Э}{{\selectfont\char221}}1 {Ю}{{\selectfont\char222}}1 {Я}{{\selectfont\char223}}1
}


\addto\captionsrussian{%
  \renewcommand{\refname}{Список использованных источников}%
}

\makeatletter
\renewcommand{\l@section}{\@dottedtocline{1}{0em}{2.3em}}
\makeatother



% ============== Заголовки разделов ==============
% Будем просто выводить их вручную жирным и по центру

% ============== Начало документа ==============
\begin{document}
\setcounter{page}{2}
% ============== Содержание и остальные разделы ==============
\tableofcontents
\newpage
\section*{Введение}
\addcontentsline{toc}{section}{Введение}

Практика студентов МГТУ им. Н.\,Э.~Баумана является обязательной частью основной образовательной программы
высшего образования и представляет собой одну из форм организации учебного процесса.

Практика — это вид учебной работы, направленный на развитие практических навыков и умений,
а также на формирование компетенций обучающихся в процессе выполнения определённых видов работ,
связанных с будущей профессиональной деятельностью.

Основными видами практики студентов Университета, обучающихся по основным образовательным программам
высшего профессионального образования, являются:
\begin{itemize}[label=--]
  \item учебная практика;
  \item производственная практика;
  \item научно-исследовательская работа (НИР).
\end{itemize}

Научно-исследовательская работа заключается в освоении студентами средств и приёмов выполнения
научно-исследовательских работ, а также в проведении собственно учебно-исследовательской деятельности.

Целями научно-исследовательской работы являются:
\begin{itemize}[label=--]
  \item овладение фундаментальной научной базой соответствующего направления подготовки,
  методологией научного творчества и современными информационными технологиями, а также подготовка
  к научно-исследовательской деятельности;
  \item подготовка материалов для выпускной квалификационной работы.
\end{itemize}

Задачами научно-исследовательской работы являются:
\begin{itemize}[label=--]
  \item участие в научно-исследовательском процессе;
  \item исследование научной темы, выданной студенту;
  \item обзор источников для выпускной квалификационной работы;
  \item выполнение теоретической части выпускной квалификационной работы;
  \item выполнение иных задач, определяемых тематикой и целями исследования.
\end{itemize}

Расчётно-пояснительная записка по результатам НИР в электронном виде (формат \texttt{Word})
после получения зачёта направляется студентом в электронный архив кафедры
по адресу: \texttt{2archive-fn@mail.ru}.

\section{Предпосылки и цели исследования}

Домен табличных данных охватывает чрезвычайно широкий круг 
прикладных задач. К нему относятся банковские и финансовые 
транзакции, результаты медицинских обследований, показатели 
промышленного мониторинга, маркетинговые профили пользователей, 
данные сенсоров кибер-физических систем и многие другие источники 
структурированной информации. Общей особенностью таких наборов 
данных является представление объектов в виде строк таблицы, 
столбцы которой являются признаками, описывающими свойства объектов 
и их взаимосвязи.

Задачи классификации табличных данных в современных прикладных областях связаны с определением класса объекта по значениям множества 
признаков. Такие данные часто содержат слабоинформативные или вовсе шумовые признаки. 
Использование исходного набора признаков без понижения их количества приводит к 
увеличению вычислительной сложности алгоритмов машинного обучения, решающих задачу классификации, и 
ухудшению их обобщающей способности. В связи с этим понижение размерности являются ключевыми 
этапами построения надёжных и эффективных моделей классификации в рамках домена табличных данных.

Понижение размерности направлено на построение компактного представления данных за счёт 
проекции исходного признакового пространства в пространство меньшей размерности, 
ориентированное прежде всего на повышение качества и устойчивости моделей классификации. 
В отличие от этого, оценка информативности признаков является алгоритмом, определяющим вклад 
каждого исходного признака в процесс принятия решения моделью, и позволяет не только выбирать 
лучшие признаки, таким образом понижая размерность, но и анализировать их значимость. 
Такой анализ особенно важен в прикладных областях, например в медицине, где требуется как 
корректно классифицировать состояние пациента, так и обосновать полученное решение, 
выявив конкретные показатели, оказывающие наибольшее влияние на постановку диагноза.

Существует большое число методов понижения размерности и оценки информативности признаков 
табличных данных, которые в общем случае можно разделить на две группы: статистические и 
нейросетевые. Статистические методы, в большинстве своем, обладают понятной интерпретацией, однако отличаются 
низкой адаптируемостью к данным со сложными взаимосвязями. Нейросетевые подходы способны 
моделировать сложные нелинейные зависимости, но требуют значительных объёмов данных, 
являются вычислительно затратными и нередко слабо обоснованы с точки зрения статистической 
надёжности и интерпретируемости.

В данной работе рассматриваются новые методы оценки информативности признаков табличных данных в 
задаче классификации, основанные на использовании унарных моделей \cite{unar}. Предлагаются два подхода: 
метод, использующий идею последовательного исключения признаков по аналогии с LOFO \cite{liu2013lofo}, и метод на 
основе ансамблей случайных подпространств, развивающий идеи RaSE \cite{tian2021rase}. Оба метода ориентированы на 
получение устойчивых и интерпретируемых оценок информативности. Эффективность предложенных 
методов исследуется на задаче бинарной классификации табличных данных и сравнивается с 
существующими решениями, как понижения размерности, так и оценки информативности.

\section{Постановка задачи}

Рассматривается задача классификации табличных данных в бинарной постановке. 
Пусть задана выборка
\begin{equation}
\mathcal{D} = \{(x_i, y_i)\}_{i=1}^n,\quad x_i \in K \subset \mathbb{R}^d,\; K \text{ - ограниченный компакт},\; y_i \in \{0,1\},
\end{equation}
где каждый объект описывается вектором признаков, включающим как информативные, 
так и слабоинформативные координаты.

Целью работы является разработка и исследование методов оценки информативности признаков 
табличных данных, направленных на выявление и отбор наиболее значимых координат признакового 
пространства для решения задачи бинарной классификации. Под оценкой информативности понимается 
построение количественных характеристик, отражающих вклад отдельных признаков в различимость 
классов и качество классификации. В рамках работы предполагается провести сравнение предложенных 
методов с существующими подходами к понижению размерности и отбору признаков.

Основное внимание уделяется разработке подходов на основе унарной классификации, подробно описанной в работе \cite{unar}, предполагающих 
построение отдельных моделей для каждого класса. На основе данной идеи предлагается разработать 
и исследовать два метода оценки информативности признаков, ориентированных на получение 
устойчивых и интерпретируемых результатов в улсовиях ограниченного объёма данных.

\section{Обзор существующих методов}

Поскольку разнообразие методов понижения размерности чрезвычайно велико, в настоящей работе
для первичного и репрезентативного сравнения выбраны, с одной стороны, классические
статистические подходы к редукции признакового пространства (метод главных компонент,
метод частичных наименьших квадратов, а также алгоритм равномерного аппроксимирования
и проекций многообразия в постановке с учителем), а с другой — достаточно мощный
нелинейный статистический метод оценки информативности признаков на основе критерия
независимости Гильберта–Шмидта с разреживающей регуляризацией. Такой выбор позволяет
сделать обзор более релевантным и сравнивать предлагаемые методы не только с
подходами отбора информативных признаков, но и с методами понижения размерности.
Ниже приведено описание, используемых методов.

\subsection{Метод главных компонент}
Метод главных компонент (англ. Principal Component Analysis) подробно описан в работах 
\cite{Pearson1901,Hotelling1933}. Он строит линейную проекцию данных в пространство меньшей 
размерности так, чтобы максимально сохранить дисперсию наблюдений. Пусть 
$X \in \mathbb{R}^{n \times d}$ - матрица центрированных наблюдений. Тогда выборочная 
ковариационная матрица имеет вид
\begin{equation}
S = \frac{1}{n} X^{\mathsf T} X,
\end{equation}
а первая главная компонента определяется как направление $\mathbf{w}_1 \in \mathbb{R}^d$, 
максимизирующее дисперсию проекций:
\begin{equation}
\mathbf{w}_1 = \arg\max_{\|\mathbf{w}\|_2=1} \mathbf{w}^{\mathsf T} S \mathbf{w}.
\end{equation}
Решение достигается выбором собственного вектора матрицы $S$, соответствующего наибольшему 
собственному значению; последующие компоненты $\mathbf{w}_2,\mathbf{w}_3,\dots$ находятся 
аналогично с дополнительным условием ортогональности к предыдущим. После выбора $q<d$ компонент 
матрица проекции $W_q=[\mathbf{w}_1,\dots,\mathbf{w}_q]$ задаёт пониженную размерность 
$Z = X W_q$.

Практически метод главных компонент полезен тем, что он прост в реализации, вычислительно 
эффективен и устраняет мультиколлинеарность признаков за счёт перехода к некоррелированным 
координатам, что нередко повышает устойчивость последующих моделей. 
В то же время метод является линейным и не требует обучения: он ориентирован на сохранение общей 
дисперсии в признаках, а не на различимость классов, поэтому направления, наиболее важные для 
классификации, могут не совпадать с построенными компонентами. Кроме того, новые координаты являются 
линейными комбинациями исходных признаков и могут быть трудны для содержательной интерпретации, 
особенно без дополнительных процедур анализа.

\subsection{Метод частичных наименьших квадратов}
Метод частичных наименьших квадратов (англ. Partial Least Squares) описан в работе 
\cite{WoldRuhe1984}. В отличие от предыдущего метода, этот метод строит латентные переменные, 
которые одновременно хорошо согласованы с меткой объекта в задаче классификации и при этом 
объясняют структуру признаков. Пусть $X \in \mathbb{R}^{n \times d}$ - центрированная матрица 
признаков, а $y \in \mathbb{R}^{n}$ - центрированный отклик. Тогда первая латентная переменная 
$t_1 = X\mathbf{w}_1$ выбирается так, чтобы ковариация с откликом была максимальной:
\begin{equation}
\mathbf{w}_1 = \arg\max_{\|\mathbf{w}\|_2=1} \mathrm{Cov}(X\mathbf{w},\,y).
\end{equation}
После нахождения $t_1$ выполняется дефляция (вычитание из $X$ и $y$ компоненты, объясняемой $t_1$), 
и процедура повторяется, формируя последовательность латентных переменных $t_1,t_2,\dots,t_q$ и 
соответствующих весов. В результате получается проекция $Z=[t_1,\dots,t_q]$, ориентированная на 
сохранение предсказательной информации относительно $y$.

С практической точки зрения метод частичных наименьших квадратов часто выигрывает у 
линейных методов, не учитывающих значения меток классов, а также повышает устойчивость к 
мультиколлинеарности и к режиму $d \gg n$. Ограничением остаётся линейный 
характер базовой постановки и необходимость выбора числа компонент, обычно по валидационной 
ошибке; при неверном выборе возможны как потеря полезной информации, так и переобучение. 
Как и в методе главных компонент, интерпретация латентных переменных требует анализа 
коэффициентов проекции, поскольку компоненты являются комбинациями исходных признаков.

\subsection{Алгоритм равномерного аппроксимирования и проекции многообразия в постановке с 
учителем}
Алгоритм равномерного аппроксимирования и проекции многообразия (англ. Uniform Manifold 
Approximation and Projection) описан в работах \cite{McInnes2018arXiv,McInnes2018JOSS}. 
Он относится к нелинейным методам и строит низкоразмерное вложение, стараясь сохранить локальную 
структуру соседства: сначала по выборке строится граф ближайших соседей, который интерпретируется 
как вероятностная структура связности, а затем подбираются точки $z_1,\dots,z_n \in \mathbb{R}^q$, 
чтобы структура связности в проекции была согласована с исходной. В базовой формулировке 
оптимизация сводится к минимизации функции потерь, эквивалентной кросс-энтропии между 
близостями $p_{ij}$ в исходном пространстве и $q_{ij}$ в проекции:
\begin{equation}
L \;=\; \sum_{i<j}\Big( -\,p_{ij}\,\log q_{ij} \;-\; (1-p_{ij})\,\log(1-q_{ij}) \Big),
\end{equation}
где $q_{ij}$ задаётся сглаженной функцией расстояния между $z_i$ и $z_j$, например семейством 
вида $q_{ij} = (1 + a\|z_i-z_j\|_2^{2b})^{-1}$ с параметрами $a,b$.

В постановке с учителем информация о метках классов используется при построении структуры 
соседства, усиливая связи внутри классов и ослабляя связи между классами, что приводит к 
проекциям, более пригодным для анализа разделимости и последующего обучения классификаторов. 
Достоинством метода является способность отражать нелинейную геометрию данных и сохранять 
локальные кластеры при понижении размерности, что особенно полезно для визуализации и 
предварительной обработки. Ограничения связаны с наличием гиперпараметров, стохастичностью 
оптимизации и слабой интерпретируемостью координат вложения: оси проекции, как правило, 
не имеют самостоятельного содержательного смысла и не допускают прямой трактовки как комбинаций 
исходных признаков.

\subsection{Метод отбора признаков на основе критерия независимости Гильберта-Шмидта и 
разреживающей регуляризации}
Критерий независимости Гильберта-Шмидта (англ. Hilbert--Schmidt Independence Criterion) подробно рассмотрен в работе \cite{Gretton2005}. 
Для выборки $\{(x_i,y_i)\}_{i=1}^n$ строятся матрицы Грама $K \in \mathbb{R}^{n\times n}$ 
и $L \in \mathbb{R}^{n\times n}$ по ядрам $k$ и $\ell$, после чего эмпирическая оценка 
зависимости выражается через центрирующую матрицу 
$H = I - \frac{1}{n}\mathbf{1}\mathbf{1}^{\mathsf T}$:
\begin{equation}
\widehat{\mathrm{HSIC}}(X,Y) = \frac{1}{(n-1)^2}\mathrm{tr}(KHLH).
\end{equation}
, где $\mathbf{1}$ - вектор из единиц длины n.

Эта форма удобна вычислительно и позволяет учитывать нелинейные зависимости за счёт выбора 
ядерных функций.

Для высокоразмерного отбора признаков используется выпуклая разреживающая постановка, предложенная, 
в частности, в работе \cite{Yamada2014}. Пусть $K_j$ --- матрица Грама, построенная только 
по $j$-й координате признакового вектора, а $\bar{K}_j = HK_jH$ и $\bar{L}=HLH$ - 
центрированные версии. Тогда веса $\beta_1,\dots,\beta_d$ находятся как решение
\begin{equation}
\min_{\beta_1,\dots,\beta_d \ge 0}\;\; \frac{1}{2}\left\|\bar{L} - \sum_{j=1}^{d}\beta_j \bar{K}_j\right\|_F^2 \;+\; \lambda\sum_{j=1}^{d}\beta_j,
\end{equation}
где $\|\cdot\|_F$ - норма Фробениуса, а $\lambda>0$ управляет разреженностью 
решения. Ненулевые коэффициенты $\beta_j$ соответствуют выбранным признакам; такая постановка 
стремится одновременно к высокой зависимости выбранных признаков с метками классов и к снижению 
избыточности между признаками, поскольку вклад каждого признака задаётся его собственной 
матрицей Грама. Сильной стороной подхода является сохранение интерпретируемости за счёт 
выбора именно исходных признаков и возможность учитывать сложные нелинейные связи через ядра, 
тогда как ограничения обычно связаны с вычислительными затратами на построение матриц Грама и 
чувствительностью к выбору ядровых гиперпараметров.

\section{Предлагаемые методы оценки информативности}

\subsection{Ограничения существующих подходов}

Так как методы оценки информативности признаков и понижения размерности в задачах классификации табличных данных чрезвычайно
разнообразны, их классификация неоднозначна, однако ее можно грубо свести к делению на 
статистические и нейросетевые подходы и обсудить проблематику каждого из семейств.

Значительная часть статистических подходов опирается на допущения о структуре зависимостей. 
В ситуациях, когда различимость классов определяется сложными нелинейными взаимодействиями 
признаков, такие методы могут давать нестабильные или неполные результаты. В свою очередь, 
нейросетевые подходы действительно способны моделировать сложные нелинейные зависимости,
но при этом возникает принципиальная трудность, связанная с корректностью вывода вне области,
поддерживаемой данными. В строгой вероятностной постановке условное математическое ожидание
\(
g(x)=\mathbb{E}(Y\mid X=x)
\)
определено единственным образом лишь почти наверное; вне носителя \(S\) распределения \(P_X\)
значения \(g(x)\) не фиксируются данными и могут быть произвольными. Следовательно,
любая модель, которая неявно продолжает правило принятия решений на весь компакт признакового
пространства, вынуждена осуществлять экстраполяцию в области \(K\setminus S\), где отсутствует 
статистическое обоснование. Это делает оценки информативности, полученные на всем компакте, 
потенциально неустойчивыми: даже небольшие изменения выборки, появление выбросов или сдвиг 
распределения, могут резко изменить поведение модели вне носителя и тем самым исказить итоговые объяснения.

Решение данной проблемы изучается в области доверенного искусственного интеллекта и требуется 
в прикладных областях, где критически важно не только получить ранжирование признаков, 
но и быть уверенным, что оно сформировано там, где вывод подкреплён данными, например, в 
медицине и инженерии. В данной работе вводится принципиальное требование: алгоритм должен иметь
механизм отказа от принятия решения, то есть формально выделять область признакового
пространства, где надёжное решение и интерпретация невозможны.

В качестве базового инструмента для реализации зоны отказа и получения статистически
обоснованных оценок далее используется идея унарной классификации с добавлением равномерного
распределения фона, предложенная в статье \cite{unar}. Такой переход позволяет построить модель,
которая обучается на наблюдаемом носителе данных и одновременно закрепляет корректное поведение
вне носителя через явную унарную постановку, что делает последующие процедуры оценки
информативности более устойчивыми и интерпретируемыми.

\subsection{Унарная классификация с фоновым классом}

Рассмотрим компакт $K=[0,1]^d$. В рамках унарной постановки классификации, для каждого класса 
строится собственный независиый классификатор: фиксируется целевой класс,
а обучающая выборка дополняется фоновыми объектами, равномерно распределёнными на $K$ и
имеющими нулевую метку. Такая постановка прямо отражает практическую идею:
модель должна уметь отличать объекты  целевого класса от любых других,
то есть заведомо не принадлежащих наблюдаемому носителю целевого распределения. На практике 
берутся $n$ наблюдений целевого класса с меткой $+1$ и добавляются независимые фоновые 
наблюдения, равномерно распределённые на $K$, с нулевой меткой. Таким образом строится по одной 
модели на каждый класс, а решение о принадлежности объекта к какому-либо из классов определяется 
на основе выходов всех унарных классиификаторов.

Ключевым объектом в такой схеме является регрессионная функция на смеси данных класса с фоном.
В \cite{unar} предлагается рассматривать модификацию исходного распределения $P$ в виде смеси
\begin{equation}
P_\alpha=(1-\alpha)P+\alpha \widetilde P,
\end{equation}
где $\widetilde P$ соответствует равномерно распределённому на $[0,1]^d$ вектору признаков с нулевой меткой.
И исследуется байесовская регрессия на модифицированном распределении
\begin{equation}
g_\alpha(x)=\mathbb{E}_{P_\alpha}(Y\mid X=x).
\end{equation}

Такая модификация позволяет зафиксировать поведение оптимального правила в местах 
не подкрепленных данными: в унарной постановке $g_\alpha(x)$ совпадает с условной вероятностью
целевого класса, а вне носителя данных обращается в нуль, формируя
естественную зону отказа.

\subsection{Теорема об экстраполяции байесовского классификатора}

В работе \cite{unar} получено строгое описание функции $g_\alpha$ и тем самым - оптимального
экстраполированного правила классификации. В частности, показывается, что решение задачи
регрессии на смеси существует, почти всюду единственно и допускает явную запись через
разложение меры $P_X$ относительно меры Лебега.
А именно, существует неотрицательная интегрируемая плотность $\rho$ и борелевское множество
$A\subseteq S$ нулевой меры Лебега, такие что для почти всех $x\in[0,1]^d$ выполнено
\begin{equation}
g_\alpha(x)=
\begin{cases}
g(x), & x\in A,\\[2mm]
\dfrac{(1-\alpha)\,g(x)\,\rho(x)}{\alpha+(1-\alpha)\rho(x)}, & \rho(x)>0,\ x\in S\setminus A,\\[3mm]
0, & \rho(x)=0,\ x\in S\setminus A\ \ \text{или}\ \ x\notin S,
\end{cases}
\end{equation}
где $g(x)=\mathbb{E}(Y\mid X=x)$.
Отсюда следует важный практический вывод: область $\{x:\,g_\alpha(x)=0\}$ покрывает
$[0,1]^d\setminus S$ почти всюду и тем самым реализует механизм отказа от решения там,
где отсутствуют данные и где стандартное обучение вынуждено делать
произвольную экстраполяцию.

Для задач оценки информативности признаков это принципиально: если модель способна
отделять область обоснованного вывода от области отказа, то важность признаков можно
определять и анализировать на доверенной области, снижая влияние выбросов,
сдвигов распределения и прочих эффектов.

\subsection{Состоятельность унарного классификатора}

Чтобы использовать унарную классификацию как основу алгоритмов информативности, для реализации доверенного метода
понижения размерности, способного улавливать сложные зависимости, необходимо понимать, 
насколько оценка $g_\alpha$ статистически надёжна и как её интерпретировать. В работе
\cite{eliseev2025} доказывается, что выход многослойного персептрона с кусочно-линейными функциями активации при обучении по
среднеквадратичной ошибке в унарной постановке может рассматриваться как состоятельная оценка гистограммной
регрессии, построенной по ячейкам линейности, индуцированным многослойным персептроном. Данный результат получен, как на практике, так и в виде следующей теоремы:

\medskip
\noindent\textbf{Теорема 1 (о сходимости нейросетевой и гистограммной регрессий) \cite{eliseev2025}.}
Пусть заданы компакт $K$, $D=\mathrm{diam}(K)$ и последовательность многослойных
персептронов с модульной функцией активации, обученных на выборке из $n$ наблюдений из распределения
с ограниченной плотностью на $K$, архитектура которых состоит из первого слоя ширины $r_n$,
$L_n$ слоёв ширины $k_n$ и одноэлементного последнего слоя с линейной активацией. Предположим, что
весовые коэффициенты инициализируются независимо из непрерывного распределения,
а параметры первого слоя заморожены при обучении. Предположим также, что целевая функция является
кусочно гладкой, $g\in \mathcal{F}_{\mathrm{MF},J,\alpha,\beta}$, и выполнены следующие ограничения:
\begin{enumerate}[label=(\roman*),leftmargin=1.25cm]
\item Число ненулевых параметров:
\begin{equation}
S_{\mathrm{nnz},n}=c_1'\max\Bigl\{n^{\frac{d}{2\beta+d}},\;n^{\frac{d-1}{\alpha+d-1}}\Bigr\}.
\end{equation}
\item Ограничение на величины параметров:
\begin{equation}
B_n\le c_1\,n^{s}.
\end{equation}
\item Количество слоёв после первого:
\begin{equation}
L_n\le c_1\Bigl(1+\max\Bigl\{\frac{\beta}{d},\;\frac{\alpha}{2(d-1)}\Bigr\}\Bigr).
\end{equation}
\item Число нейронов в первом слое:
\begin{equation}
r_n>2d.
\end{equation}
\item Ограничение на скорость роста архитектуры:
\begin{equation}
k_n^{\,L_n\min\{d,k_n\}}\,r_n^{d}=o(n),\qquad
\frac{k_n^{\,L_n\min\{d,k_n\}}\,r_n^{d}\;d(k_nL_n+r_n)\log(k_nL_n+r_n)\,\log n}{n}\to 0.
\end{equation}
\item Ограничения на ширину первого слоя (для произвольного $\gamma>0$):
\begin{equation}
\sum_{n=1}^{\infty}\exp(-c\gamma^{4}r_n)<\infty,\qquad
r_n> C\gamma^{-12}\frac{\omega(K)^2}{D^{10}}.
\end{equation}
\end{enumerate}
Тогда
\begin{equation}
c_n(X)-h_n(X)\xrightarrow{P}0,\qquad n\to\infty.
\end{equation}

Привеленные теоремы из работ по унарной классификации дают два ключевых вывода, которые становятся фундаментом для построения
доверенных оценок информативности признаков. Во-первых, унарная нейросетевая модель,
обученная на смеси распределения целевого класса и распределния равномерного фона, имеет строгую статистическую опору:
её выход можно рассматривать как корректную асимптотическую оценку байесовской регрессии
$g_\alpha$, а значит, она наследует принципиальное свойство унарной постановки - формирование
естественной зоны отказа вне носителя данных.
Во-вторых, тот же результат объясняет интерпретируемость: нейросетевая регрессия
асимптотически эквивалентна гистограммной регрессии $h_n$ на ячейках линейности. В совокупности это означает, что
унарная классификация одновременно сочетает нелинейную выразительность с математически
обоснованным поведением в областях, где нет данных, и интерпретируемой структурой оценки, что делает
её удобной базой для конструирования устойчивых методов оценки информативности признаков.

Далее будут предложены две новые идеи оценки информативности 
признаков табличных данных с использованием унарных классификаторов.

\subsection{Метод понижения размерности на основе LOFO}

\subsubsection{Идея метода и постановка}
В качестве первого предлагаемого метода рассматривается подход оценки информативности признаков,
построенный по схеме последовательного исключения признаков по одному с переобучением моделей после исключения каждого признака (Leave-One-Feature-Out).
Схема LOFO подробно описана в статье \cite{liu2013lofo} и в данной работе 
адаптирована к классификации в унарной постоновке.
Ключевая цель метода сохранить подмножество признаков фиксированного размера $q<d$,
которое обеспечивает максимальную различимость классов в смысле унарных метрик, определённых ниже.

\subsubsection{Унарные модели и зона отказа}
Пусть в бинарной задаче выделены два класса $\omega_1$ и $\omega_2$, и для каждого класса обучается
отдельная унарная модель, отличающая объекты данного класса от фона.
Обозначим через $c_1(x)$ и $c_2(x)$ выходы двух унарных классификаторов, построенных для классов
$\omega_1$ и $\omega_2$ соответственно. В соответствии с предыдущими разделами, данные каждого класса
дополняются фоновыми объектами, равномерно распределёнными на компакте $K$, что формирует
статистически обоснованную зону отказа: вне носителя данных унарные выходы стремятся к нулю,
а решение становится консервативным.

\subsubsection{Пороговое разбиение пространства}
Для перехода от вещественных выходов $c_1(x),c_2(x)$ к определению класса вводятся пороги
$\beta_1,\beta_2\in\mathbb{R}$. На их основе пространство признаков разбивается на четыре области:
\begin{align}
R_1(\beta_1,\beta_2) &= \{x:\; c_1(x)\ge \beta_1,\; c_2(x)<\beta_2\},\\
R_2(\beta_1,\beta_2) &= \{x:\; c_2(x)\ge \beta_2,\; c_1(x)<\beta_1\},\\
R_{12}(\beta_1,\beta_2) &= \{x:\; c_1(x)\ge \beta_1,\; c_2(x)\ge \beta_2\},\\
R_{0}(\beta_1,\beta_2) &= \{x:\; c_1(x)< \beta_1,\; c_2(x)< \beta_2\}.
\end{align}
Здесь $R_1$ и $R_2$ соответствуют однозначному принятию решения в пользу одного из классов,
$R_{12}$ описывает область наложения,
а $R_0$ соответствует зоне отказа, что является принципиальным отличием
от стандартной бинарной схемы без механизма отказа.

\subsubsection{Подбор порогов по критерию $J_k$}
Пороги $\beta_1,\beta_2$ должны удовлетворять двум требованиям: высокая доля корректно распознанных
объектов целевых классов и низкая доля ложных срабатываний на фоне, а также
на противоположном классе. Для каждого класса $k\in\{1,2\}$ можно рассматривать критерий в виде:
\begin{equation}
J_k(\beta_k) = \mathrm{TPR}_k(\beta_k) + \mathrm{TNR}_k(\beta_k) - 1,
\end{equation}
где $\mathrm{TPR}_k(\beta_k)$ есть доля объектов класса $\omega_k$ с $c_k(x)\ge \beta_k$,
а $\mathrm{TNR}_k(\beta_k)$ есть доля объектов фона и, объектов второго класса,
для которых выполняется $c_k(x)<\beta_k$. Максимизация $J_k(\beta_k)$ по $\beta_k$ даёт
порог, который балансирует чувствительность и специфичность унарного классификатора
в присутствии фона.

\subsubsection{Счётчики попаданий и доли $a_1,a_2,b_1,b_2$}
Пусть на валидационной выборке заданы подвыборки классов
$\mathcal{D}^{(1)}_{\mathrm{val}}=\{x_i:\; y_i=\omega_1\}$ и
$\mathcal{D}^{(2)}_{\mathrm{val}}=\{x_i:\; y_i=\omega_2\}$ размеров $n_1$ и $n_2$.
Введём счётчики попаданий объектов каждого класса в области $R_1,R_2,R_{12}$:
\begin{align}
n_{10}^{(1)} &= \sum_{i:\,y_i=\omega_1} \mathbf{1}\{x_i\in R_1(\beta_1,\beta_2)\},\\
n_{02}^{(2)} &= \sum_{i:\,y_i=\omega_2} \mathbf{1}\{x_i\in R_2(\beta_1,\beta_2)\},\\
n_{12}^{(1)} &= \sum_{i:\,y_i=\omega_1} \mathbf{1}\{x_i\in R_{12}(\beta_1,\beta_2)\},\\
n_{12}^{(2)} &= \sum_{i:\,y_i=\omega_2} \mathbf{1}\{x_i\in R_{12}(\beta_1,\beta_2)\}.
\end{align}
Далее определим доли однозначного распознавания и доли наложения:
\begin{equation}
a_1=\frac{n_{10}^{(1)}}{n_1},\qquad a_2=\frac{n_{02}^{(2)}}{n_2},\qquad
b_1=\frac{n_{12}^{(1)}}{n_1},\qquad b_2=\frac{n_{12}^{(2)}}{n_2}.
\end{equation}
Величины $a_1,a_2$ измеряют, какая доля объектов каждого класса попадает в область
однозначного решения, то есть характеризуют эффективность разделения с учётом зоны отказа.
Величины $b_1,b_2$ измеряют долю объектов каждого класса, попадающих в область наложения,
то есть характеризуют неоднозначность решения.

\subsubsection{Унарные метрики $F_{12}$, $G_{12}$ и агрегированный критерий $S$}
Для получения симметричных по классам скалярных характеристик используются гармонические средние:
\begin{equation}
F_{12}=\frac{2a_1a_2}{a_1+a_2+\varepsilon},\qquad
G_{12}=\frac{2b_1b_2}{b_1+b_2+\varepsilon},
\end{equation}
где $\varepsilon>0$ - малый параметр, предотвращающий деление на нуль.
Метрика $F_{12}$ является аналогом $F1$-меры, отражая согласованно высокую эффективность
по обоим классам, тогда как $G_{12}$ отражает степень наложения классов, которая нежелательна.
Для окончательной оценки качества вводится агрегированный критерий
\begin{equation}
S(\beta_1,\beta_2;c_G)=F_{12}-c_G\,G_{12},\qquad c_G>0,
\end{equation}
где параметр $c_G$ задаёт штраф за наложение. Большие значения $S$ соответствуют ситуациям,
когда классы разделяются эффективно и при этом редко попадают в область конфликта $R_{12}$.

\subsubsection{Итеративное исключение признаков}
Пусть $J\subseteq\{1,\dots,d\}$ - текущий набор сохраняемых признаков.
Для каждого кандидата $j\in J$ строится редуцированная выборка, в которой удалена координата $j$,
то есть рассматривается проекция $x_{-j}\in\mathbb{R}^{|J|-1}$. На этой редуцированной выборке
переобучаются унарные модели $c_{1,-j}$ и $c_{2,-j}$, выбираются пороги $\beta_{1,-j},\beta_{2,-j}$
и вычисляется значение критерия $S_{-j}$ на валидации.
После этого исключается признак
\begin{equation}
j^\star = \arg\max_{j\in J} S_{-j},
\end{equation}
то есть удаляется координата, исключение которой приводит к наибольшему улучшению или наименьшему ухудшению разделимости в смысле унарных метрик. Процедура повторяется итеративно до тех пор,
пока не останется $q$ признаков. Итогом работы метода является подмножество признаков $J$,
которое можно использовать как результат понижения размерности.

Код реализации предлагаемого метода представлен в приложении А.

\subsubsection{Вычислительная сложность и ограничения}
Схема LOFO требует переобучения моделей при каждом исключении признака.
В простейшей реализации вычислительная сложность растёт как $O(d(d-q))$ обучений унарных моделей,
что ограничивает применение при больших $d$.

Дополнительным ограничением является проклятие размерности, связанное с использованием
равномерного фонового распределения в процессе обучения моделей. В высоких размерностях объём гиперкуба быстро растёт,
и количество данных фона, необходимых для полноценного покрытия всего гиперкуба, растет экспоненциально.

\subsection{Метод понижения размерности на основе ансамблей случайных подпространств и SHAP-оценок}

В качестве второго предлагаемого метода рассматривается подход, сочетающий идею ансамблей
случайных подпространств и унарную классификацию с зоной отказа, а также использование
аддитивных объяснений Шепли для обновления оценок информативности признаков.
Цель метода - получить устойчивый вектор информативности исходных координат и затем выбрать
$q<d$ наиболее значимых признаков, при этом поборов проклятие размерности, наследуемое от обучения унарных классификаторов.

\subsubsection{Ансамбли случайных подпространств (RaSE)}

Метод ансамблей случайных подпространств (англ.\ Random Subspace Ensemble), известный как RaSE,
предлагается в работе \cite{tian2021rase}. Его основная идея состоит в том, что вместо обучения
одной сложной модели на всём пространстве $\mathbb{R}^d$ строится ансамбль слабых моделей,
каждая из которых обучается на подпространстве признаков небольшой размерности. Подпространства
выбираются случайно, но не равномерно: вероятность включения координат адаптивно обновляется так,
чтобы чаще выбирать признаки, встречающиеся в наиболее качественных подпространствах.
В \cite{tian2021rase} приведён детальный теоретический анализ и обоснование такой схемы.

Достоинство данной идеи в контексте унарной классификации заключается в следующем:
при фиксированном размере подпространства $m\ll d$ обучение и валидация унарных моделей
становятся существенно устойчивее и дешевле, а генерация фоновых объектов равномерного шума
в $m$-мерном гиперркубе требует существенно меньшего числа точек, чем в $d$-мерном.

\subsubsection{Аддитивные объяснения Шепли}

Для измерения вклада координат внутри каждого выбранного подпространства используются SHAP-оценки
(англ.\ SHapley Additive explanations), предложенные в работе \cite{lundberg2017shap}.
SHAP рассматривает объяснение предсказания как аддитивную модель по признакам:
\begin{equation}
f(x) \;=\; \phi_0 \;+\; \sum_{j=1}^{d} \phi_j(x),
\end{equation}
где $f(x)$ - выход модели в точке $x$, $\phi_0$ - базовый уровень,
а $\phi_j(x)$ - вклад $j$-й координаты.

Вклад признака $\phi_j(x)$ определяется как значение Шепли из теории игр:
\begin{equation}
\phi_j(x)
\;=\;
\sum_{S \subseteq \{1,\dots,d\}\setminus\{j\}}
\frac{|S|!\,(d-|S|-1)!}{d!}
\Bigl(
f_{S\cup\{j\}}(x_{S\cup\{j\}}) - f_S(x_S)
\Bigr),
\end{equation}
где $S$ - подмножество признаков, не содержащее $j$,
$f_S(x_S)$ - значение модели, вычисленное при использовании только признаков из $S$,
а остальные координаты считаются отсутствующими.

В \cite{lundberg2017shap} показано, что в классе аддитивных объяснений, удовлетворяющих
аксиомам локальной точности, симметрии и отсутствия вклада для неиспользуемых признаков,
существует единственное решение, совпадающее со значениями Шепли,
что делает SHAP принципиально обоснованной мерой важности признаков.

\subsubsection{Описание предлагаемого алгоритма RaSE-SHAP в унарной постановке}

Пусть решается бинарная задача с классами $\omega_1,\omega_2$. Как и ранее, обучаются два унарных
классификатора: $c_1(x)$ для класса $\omega_1$ и $c_2(x)$ для класса $\omega_2$, причём обучение
каждого унарного классификатора проводится на множестве целевых объектов, смешанных с фоновыми объектами,
равномерно распределёнными на компакте в соответствующем подпространстве признаков.

Метод конструирует вектор вероятностей выбора координат
\begin{equation}
\pi=(\pi_1,\dots,\pi_d),\qquad \pi_j\ge 0,\quad \sum_{j=1}^d \pi_j=1,
\end{equation}
который интерпретируется как текущая оценка информативности признаков. Итерации обновления веектора вероятностей состоят
из трёх этапов.

\medskip
\noindent\textbf{Генерация и отбор подпространств.}
Фиксируется размер подпространства $m$ и число кандидатов. На каждой итерации генерируются
подпространства-кандидаты $J\subset\{1,\dots,d\}$, $|J|=m$:
с вероятностью $\lambda$ координаты выбираются равномерно, а с вероятностью $1-\lambda$
- согласно распределению $\pi$, что обеспечивает возможность преодолеть проблемы зависания метода.
Для каждого кандидата $J$ обучаются унарные модели $c_{1,J}$ и $c_{2,J}$ на признаках $x_J$ и
на валидации подбираются пороги $\beta_{1,J},\beta_{2,J}$, максимизирующие критерий разделимости
\begin{equation}
S_J = F_{12}(\beta_{1,J},\beta_{2,J}) - c_G\,G_{12}(\beta_{1,J},\beta_{2,J}),
\end{equation}
где $F_{12}$ и $G_{12}$ определены согласно предыдущему методу, а $c_G>0$ задаёт штраф за область наложения.
Далее сохраняется набор лучших подпространств и соответствующих моделей по величине $S_J$.

\medskip
\noindent\textbf{Оценка важности координат внутри лучших подпространств через SHAP.}
Для каждого выбранного лучшего подпространства $J$ вычисляются SHAP-оценки для унарных моделей.
Обозначим через $\phi^{(1)}_{J}(x)\in\mathbb{R}^{m}$ вектор SHAP-вкладов для модели $c_{1,J}$,
а через $\phi^{(2)}_{J}(x)$ - для модели $c_{2,J}$. Для устойчивости далее используются только
неотрицательные вклады, так как только они влияют положительно на принадлежность данных к конкретному классу, и усредним результат по объектам соответствующего класса на валидации:
\begin{equation}
u^{(1)}_{J}=\frac{1}{|V_1|}\sum_{x\in V_1}\max(\phi^{(1)}_{J}(x),0),\qquad
u^{(2)}_{J}=\frac{1}{|V_2|}\sum_{x\in V_2}\max(\phi^{(2)}_{J}(x),0),
\end{equation}
где $V_1$ и $V_2$ - валидационные подвыборки классов $\omega_1$ и $\omega_2$.
Затем векторы нормируются на сумму и объединяются в симметричную оценку значимости координат
подпространства через гармоническое среднее:
\begin{equation}
v_J=\frac{2\,u^{(1)}_{J}\odot u^{(2)}_{J}}{u^{(1)}_{J}+u^{(2)}_{J}+\varepsilon},
\end{equation}
где $\odot$ - поэлементное умножение, $\varepsilon>0$ - малый параметр.
Такое объединение усиливает признаки, которые одновременно важны для обоих унарных классификаторов,
и подавляет признаки, значимые лишь для одного из них.

\medskip
\noindent\textbf{Обновление распределения $\pi$.}
Пусть $\mathcal{J}_{\mathrm{top}}$ - набор лучших подпространств на текущей итерации.
Из векторов $v_J$ формируется общий список информативности по всем признакам согласно отобранным кандидатам подпространств.
\begin{equation}
w_i=\sum_{J\in\mathcal{J}_{\mathrm{top}}}\mathbf{1}\{i\in J\}\, (v_J)_{\mathrm{pos}(i,J)},
\end{equation}
где $\mathrm{pos}(i,J)$ - позиция признака $i$ внутри упорядоченного списка индексов $J$.
После нормировки $w$ по полученным результатам с использованием сглаживания обновляется общий вектор информативности признаков
\begin{equation}
\pi^{\mathrm{new}}=(1-\alpha)\pi+\alpha\,\frac{w}{\sum_{j=1}^d w_j},
\qquad \alpha\in(0,1),
\end{equation}
а затем выполняется нормировка с ограничениями снизу и сверху
$L\le \pi_j\le U$, предотвращающая вырождение распределения.
\begin{equation}
\pi \leftarrow \arg\min_{p}\|p-\pi^{\mathrm{new}}\|_2^2
\quad \text{при}\quad
\sum_{j=1}^d p_j=1,\;\; L\le p_j\le U.
\end{equation}
Наличие нижней границы $L$ обеспечивает неизчезающую вероятность исследования координат, а верхняя
граница $U$ не позволяет $\pi$ слишком рано схлопнуться в несколько признаков.

После выполнения заданного числа итераций итоговая оценка информативности задаётся вектором $\pi$.
Как результат понижения размерности выбирается множество индексов $q$ наибольших компонентов $\pi$.
Дополнительно помимо вектора информативности признаков в качестве результата работы метода получен ансамбль унарных
классифиаторов, который решает задачу классификации в исходном пространстве признаков.

Код реализации данного метода находится в приложении Б.

\subsubsection{Преимущества и ограничения метода}

Ключевым преимуществом предлагаемого метода является частичное преодоление проклятия размерности.
В отличие от прямого обучения унарных моделей на всём пространстве признаков, обучение и валидация
выполняются на подпространствах существенно меньшей размерности, что снижает требования к числу
наблюдений и к объёму фоновых данных, необходимых для формирования корректной зоны отказа.
Адаптивный характер выбора подпространств позволяет концентрировать вычислительные ресурсы на
наиболее информативных координатах и тем самым получать более устойчивые оценки информативности
в условиях высокой размерности и ограниченного объёма выборки.

В то же время метод остаётся вычислительно сложным. На каждой итерации требуется обучение
нескольких унарных моделей на различных подпространствах, подбор порогов и вычисление SHAP-оценок,
что приводит к значительным затратам времени и памяти, особенно при большом числе признаков
и при использовании нейросетевых базовых моделей. Таким образом, практическое применение метода
требует компромисса между точностью оценок информативности и доступными вычислительными ресурсами.

Дополнительным ограничением является необходимость задания набора гиперпараметров и констант,
включая размер подпространств, число моделей в ансамбле, параметры сглаживания распределения
вероятностей координат и границы допустимых значений этих вероятностей. Некорректный выбор данных
констант может приводить как к излишне медленной сходимости алгоритма, так и к преждевременному
вырождению распределения информативности. Вопрос о теоретически обоснованном и практически
устойчивом выборе таких параметров в различных режимах работы алгоритма требует отдельного
исследования и будет рассматриваться в дальнейшем на основе более глубоких теоретических
рассуждений и анализа свойств предлагаемого метода.

\section{Метрики оценки качества понижения размерности}

Оценка качества понижения размерности в данной работе проводится по тому, как редуцированное пространство
влияет на обучение и качество произвольных моделей классификации.
Такой подход является принципиальным: цель понижения размерности в задачах классификации
заключается не в сохранении структуры данных как таковой, а в сохранении или улучшении
предсказательной способности моделей, обучаемых на уменьшенном наборе признаков.

В экспериментальном исследовании для этой цели используются две модели:
многослойный персептрон и градиентный бустинг над гистограммной регрессией.
Для каждой модели проводится обучение на исходном пространстве признаков и на пространстве,
полученном после применения метода понижения размерности до фиксированной размерности $q$.
Качество сравнивается по стандартным метрикам бинарной классификации,
включая площадь под ROC-кривой (ROC AUC) и площадь под PR-кривой (PR AUC).

Для вывода модели $s(x)\in\mathbb{R}$, по которому строится бинарное решение
$\widehat y_\tau=\mathbf{1}\{s(x)\ge \tau\}$ при пороге $\tau$,
обозначим через $\mathrm{TP}(\tau),\mathrm{FP}(\tau),\mathrm{TN}(\tau),\mathrm{FN}(\tau)$
числа истинно положительных, ложно положительных, истинно отрицательных и ложно отрицательных
предсказаний при пороге $\tau$. Тогда вводятся
\begin{equation}
\mathrm{TPR}(\tau)=\frac{\mathrm{TP}(\tau)}{\mathrm{TP}(\tau)+\mathrm{FN}(\tau)},\qquad
\mathrm{FPR}(\tau)=\frac{\mathrm{FP}(\tau)}{\mathrm{FP}(\tau)+\mathrm{TN}(\tau)}.
\end{equation}
ROC-кривая - это зависимость $\mathrm{TPR}(\tau)$ от $\mathrm{FPR}(\tau)$ при изменении $\tau$,
а метрика ROC AUC определяется как площадь под этой кривой:
\begin{equation}
\mathrm{ROC\ AUC}=\int_{0}^{1}\mathrm{TPR}(u)\,du,
\qquad u=\mathrm{FPR}(\tau).
\end{equation}
Эквивалентная вероятностная интерпретация состоит в том, что ROC AUC равна вероятности того,
что случайно выбранный объект положительного класса получит больший вывод модели, чем вывод случайно выбранного
объекта отрицательного класса:
\begin{equation}
\mathrm{ROC\ AUC}=\mathbb{P}\bigl(s(X^+)>s(X^-)\bigr)+\frac{1}{2}\mathbb{P}\bigl(s(X^+)=s(X^-)\bigr).
\end{equation}

Для PR-кривой используются точность (precision) и полнота (recall). При пороге $\tau$:
\begin{equation}
\mathrm{Precision}(\tau)=\frac{\mathrm{TP}(\tau)}{\mathrm{TP}(\tau)+\mathrm{FP}(\tau)},\qquad
\mathrm{Recall}(\tau)=\mathrm{TPR}(\tau)=\frac{\mathrm{TP}(\tau)}{\mathrm{TP}(\tau)+\mathrm{FN}(\tau)}.
\end{equation}
PR-кривая - это зависимость $\mathrm{Precision}(\tau)$ от $\mathrm{Recall}(\tau)$ при изменении $\tau$,
а метрика PR AUC (средняя точность) определяется как площадь под PR-кривой:
\begin{equation}
\mathrm{PR\ AUC}=\int_{0}^{1}\mathrm{Precision}(r)\,dr,
\qquad r=\mathrm{Recall}(\tau).
\end{equation}
Данная метрика особенно чувствительна к качеству ранжирования в области малых ложноположительных
срабатываний и обычно более информативна при дисбалансе классов.

Дополнительно анализируются разности
\[
\Delta \mathrm{AUC} = \mathrm{AUC}_{\mathrm{red}} - \mathrm{AUC}_{\mathrm{base}}, \qquad
\Delta \mathrm{PR} = \mathrm{PR}_{\mathrm{red}} - \mathrm{PR}_{\mathrm{base}},
\]
которые напрямую отражают выигрыш или потерю качества вследствие понижения размерности
для фиксированной модели.

Помимо предсказательных характеристик, учитываются вычислительные затраты.
Измеряется время работы алгоритма понижения размерности $t_{\mathrm{reducer}}$,
а также время обучения модели на редуцированном и на исходном пространствах признаков
$t_{\mathrm{model}}$ и $t_{\mathrm{model,base}}$ соответственно.
Такой набор метрик позволяет оценивать методы понижения размерности комплексно,
с учётом компромисса между качеством классификации и вычислительной эффективностью,
что особенно важно в прикладных сценариях с ограниченными вычислительными ресурсами.

% ======================================================================
%  Экспериментальное исследование (4 датасета, 6 таблиц на датасет)
% ======================================================================

\section{Экспериментальное исследование}

В экспериментальном исследовании сравниваются два предложенных метода
- унарный LOFO (MLP\_unar) и ансамбль случайных подпространств с SHAP-оценками (Ensembleunar) -
с базовыми методами из обзора: метод главных компонент (PCA), метод частичных наименьших квадратов (PLS),
алгоритм Uniform Manifold Approximation and Projection в постановке с учителем (UMAP\_sup),
а также метод отбора признаков на основе критерия независимости Гильберта-Шмидта и разреживающей регуляризации (HSIC\_Lasso).
Отметим, что в сравнении присутствуют как методы понижения размерности (PCA, PLS, UMAP\_sup),
так и методы оценки информативности (HSIC\_Lasso, MLP\_unar, Ensembleunar),
по этой причине методы оценки информативности, только отбирающие существующие признаки, не всегда могут лидировать.
однако в данной работе рассматриваются оба семейства методов для более полного анализа.
В дальнейшем планируется расширить сравнение за счёт большего числа альтернативных методов и стратегий настройки.

Способы оценки качества работы модели, используемые в экспериментах, описаны ранее.

Код реализации всех экспериментов находится в приложении В.

\captionsetup{labelsep=period,justification=centering,font=small}
\renewcommand{\arraystretch}{1.15}

% ----------------------------------------------------------------------
\subsection{Синтетический набор данных DS-A}

Набор данных DS-A представляет собой синтетическую бинарную задачу классификации из $n=3000$ объектов,
описанных $d=8$ признаками. При генерации данных задаётся, что $3$ признака являются информативными,
ещё $2$ - избыточными, то есть линейно связаны с информативными и не добавляют
существенно новой информации. Классы разделены с параметром $\texttt{class\_sep}=1.2$, а в разметку
внесён небольшой уровень шума: доля случайных инверсий меток составляет $\texttt{flip\_y}=0.02$.
Чтобы сделать зависимость целевого признака от части координат нелинейной, дополнительно выполняются
преобразования признаков: первая координата заменяется на $x_0\leftarrow \sin(x_0)+0.1x_1$,
а третья - на $x_2\leftarrow x_2^2$. В результате DS-A сочетает линейно-информативную структуру,
избыточность и локальные нелинейности, что позволяет оценить устойчивость методов редукции и отбора
признаков при уменьшении размерности. В экспериментах для DS-A рассматривались размеры редукции
$q\in\{2,3,5\}$.

% ---- DS-A, HGBT ----
\begin{table}[H]
\caption{Результаты на DS-A (HGBT, $q=2$)}
\label{tab:exp:DSA:HGBT:q2}
\centering\small
\resizebox{\textwidth}{!}{%
\begin{tabular}{lrrrrrr}
\hline
Метод & ROC AUC & PR AUC & $\Delta$ROC AUC & $\Delta$PR AUC & $t_{\text{ред}}$, c & $t_{\text{мод}}$, c\\
\hline
HSIC\_Lasso & 0.970565 & 0.974506 & -0.017788 & -0.014287 & 0.285843 & 0.200469\\
MLP\_unar   & 0.986348 & 0.986826 & -0.002006 & -0.001966 & 117.725034 & 0.204162\\
PCA         & 0.978915 & 0.975961 & -0.009438 & -0.012831 & 0.001453 & 0.199110\\
PLS         & 0.980662 & 0.979784 & -0.007691 & -0.009009 & 0.001888 & 0.199345\\
UMAP\_sup   & 0.980732 & 0.980909 & -0.007621 & -0.007884 & 11.710310 & 0.199346\\
\hline
\end{tabular}}
\end{table}

\begin{table}[H]
\caption{Результаты на DS-A (HGBT, $q=3$)}
\label{tab:exp:DSA:HGBT:q3}
\centering\small
\resizebox{\textwidth}{!}{%
\begin{tabular}{lrrrrrr}
\hline
Метод & ROC AUC & PR AUC & $\Delta$ROC AUC & $\Delta$PR AUC & $t_{\text{ред}}$, c & $t_{\text{мод}}$, c\\
\hline
HSIC\_Lasso & 0.987754 & 0.988550 & -0.000599 & -0.000243 & 0.300731 & 0.211273\\
MLP\_unar   & 0.985921 & 0.987263 & -0.002432 & -0.001529 & 106.325773 & 0.323044\\
PCA         & 0.982696 & 0.981152 & -0.005657 & -0.007641 & 0.001158 & 0.210256\\
PLS         & 0.981381 & 0.981300 & -0.006972 & -0.007492 & 0.002058 & 0.206862\\
UMAP\_sup   & 0.981740 & 0.978744 & -0.006613 & -0.010049 & 12.419136 & 0.215452\\
\hline
\end{tabular}}
\end{table}

\begin{table}[H]
\caption{Результаты на DS-A (HGBT, $q=5$)}
\label{tab:exp:DSA:HGBT:q5}
\centering\small
\resizebox{\textwidth}{!}{%
\begin{tabular}{lrrrrrr}
\hline
Метод & ROC AUC & PR AUC & $\Delta$ROC AUC & $\Delta$PR AUC & $t_{\text{ред}}$, c & $t_{\text{мод}}$, c\\
\hline
HSIC\_Lasso & 0.986532 & 0.987296 & -0.001821 & -0.001497 & 0.280258 & 0.223886\\
MLP\_unar   & 0.988919 & 0.990394 & 0.000566 & 0.001601 & 75.101636 & 0.231676\\
PCA         & 0.986327 & 0.983510 & -0.002027 & -0.005282 & 0.001225 & 0.236653\\
PLS         & 0.983284 & 0.983515 & -0.005069 & -0.005277 & 0.002653 & 0.226112\\
UMAP\_sup   & 0.981842 & 0.975839 & -0.006511 & -0.012954 & 12.897769 & 0.223852\\
\hline
\end{tabular}}
\end{table}

% ---- DS-A, MLP ----
\begin{table}[H]
\caption{Результаты на DS-A (MLP, $q=2$)}
\label{tab:exp:DSA:MLP:q2}
\centering\small
\resizebox{\textwidth}{!}{%
\begin{tabular}{lrrrrrr}
\hline
Метод & ROC AUC & PR AUC & $\Delta$ROC AUC & $\Delta$PR AUC & $t_{\text{ред}}$, c & $t_{\text{мод}}$, c\\
\hline
HSIC\_Lasso & 0.971751 & 0.975836 & -0.014325 & -0.010527 & 0.285843 & 0.137024\\
MLP\_unar   & 0.984589 & 0.985793 & -0.001487 & -0.000570 & 117.725034 & 0.205935\\
PCA         & 0.984104 & 0.984548 & -0.001971 & -0.001815 & 0.001453 & 0.261257\\
PLS         & 0.984433 & 0.984184 & -0.001642 & -0.002179 & 0.001888 & 0.156200\\
UMAP\_sup   & 0.977611 & 0.975042 & -0.008465 & -0.011321 & 11.710310 & 0.147092\\
\hline
\end{tabular}}
\end{table}

\begin{table}[H]
\caption{Результаты на DS-A (MLP, $q=3$)}
\label{tab:exp:DSA:MLP:q3}
\centering\small
\resizebox{\textwidth}{!}{%
\begin{tabular}{lrrrrrr}
\hline
Метод & ROC AUC & PR AUC & $\Delta$ROC AUC & $\Delta$PR AUC & $t_{\text{ред}}$, c & $t_{\text{мод}}$, c\\
\hline
HSIC\_Lasso & 0.983300 & 0.984784 & -0.002776 & -0.001579 & 0.300731 & 0.273864\\
MLP\_unar   & 0.984815 & 0.986107 & -0.001260 & -0.000256 & 106.325773 & 0.827545\\
PCA         & 0.984529 & 0.984277 & -0.001547 & -0.002086 & 0.001158 & 0.191964\\
PLS         & 0.984493 & 0.984251 & -0.001582 & -0.002112 & 0.002058 & 0.157493\\
UMAP\_sup   & 0.979955 & 0.979153 & -0.006120 & -0.007210 & 12.419136 & 0.152262\\
\hline
\end{tabular}}
\end{table}

\begin{table}[H]
\caption{Результаты на DS-A (MLP, $q=5$)}
\label{tab:exp:DSA:MLP:q5}
\centering\small
\resizebox{\textwidth}{!}{%
\begin{tabular}{lrrrrrr}
\hline
Метод & ROC AUC & PR AUC & $\Delta$ROC AUC & $\Delta$PR AUC & $t_{\text{ред}}$, c & $t_{\text{мод}}$, c\\
\hline
HSIC\_Lasso & 0.986713 & 0.988009 & 0.000638 & 0.001646 & 0.280258 & 0.173197\\
MLP\_unar   & 0.981584 & 0.982671 & -0.004491 & -0.003692 & 75.101636 & 0.188049\\
PCA         & 0.984573 & 0.983518 & -0.001502 & -0.002845 & 0.001225 & 0.255150\\
PLS         & 0.982666 & 0.983012 & -0.003409 & -0.003351 & 0.002653 & 0.166903\\
UMAP\_sup   & 0.980662 & 0.979699 & -0.005413 & -0.006664 & 12.897769 & 0.209122\\
\hline
\end{tabular}}
\end{table}

На DS-A метод MLP\_unar показывает наилучшее качество для HGBT при $q=5$ и для MLP при $q\in\{2,3\}$,
однако различия с ближайшими конкурентами в этих режимах находятся на уровне тысячных.
При $q=5$ и базовой модели MLP метод уступает HSIC\_Lasso и линейным редукторам, однако разница по ROC AUC опять же порядка $0.005$ и
остаётся в пределах сопоставимого качества; такое поведение может зависеть от статистической корректности оценки,
а также от настройки гиперпараметров унарной части.

% ----------------------------------------------------------------------
\subsection{Синтетический набор данных DS-B}

Набор данных DS-B представляет собой синтетическую бинарную задачу классификации повышенной размерности,
содержащую $n=6000$ объектов, описанных $d=30$ признаками. При генерации данных задаётся,
что $5$ признаков являются информативными, ещё $5$ - избыточными, а распределение классов
существенно несбалансировано: доля объектов положительного класса составляет около $25\%$.
Разделимость классов умеренная ($\texttt{class\_sep}=0.8$), при этом в разметку добавлен небольшой
уровень шума с долей инверсий меток $\texttt{flip\_y}=0.01$.

Для усложнения структуры зависимости целевого признака от признакового пространства в набор данных
включены как нелинейные преобразования, так и взаимодействия между координатами. В частности,
вводятся признаки вида $x_{10}\leftarrow \sin(x_0)\,x_1$, $x_{11}\leftarrow x_2^2$ и
$x_{12}\leftarrow \exp(-x_3^2)$, а также набор слабо зашумлённых копий первых информативных признаков,
размещённых в координатах $20{:}25$. Такая конструкция моделирует типичную для прикладных табличных данных
ситуацию, в которой одновременно присутствуют информативные и избыточные признаки, коррелированные группы
переменных и смешанная линейно-нелинейная структура зависимостей.

В экспериментах для DS-B рассматривались размеры редукции $q\in\{5,10,15\}$.

% ---- DS-B, HGBT ----
\begin{table}[H]
\caption{Результаты на DS-B (HGBT, $q=5$)}
\label{tab:exp:DSB:HGBT:q5}
\centering\small
\resizebox{\textwidth}{!}{%
\begin{tabular}{lrrrrrr}
\hline
Метод & ROC AUC & PR AUC & $\Delta$ROC AUC & $\Delta$PR AUC & $t_{\text{ред}}$, c & $t_{\text{мод}}$, c\\
\hline
Ensembleunar & 0.943656 & 0.872110 & -0.014764 & -0.033804 & 1462.708774 & 0.338188\\
HSIC\_Lasso  & 0.954001 & 0.896141 & -0.004419 & -0.009773 & 0.605554 & 0.337242\\
PLS          & 0.946973 & 0.885809 & -0.011447 & -0.020106 & 0.007177 & 0.297453\\
PCA          & 0.677136 & 0.429959 & -0.281284 & -0.475956 & 0.002488 & 0.285192\\
UMAP\_sup    & 0.745107 & 0.527121 & -0.213313 & -0.378794 & 16.514097 & 0.295961\\
\hline
\end{tabular}}
\end{table}

\begin{table}[H]
\caption{Результаты на DS-B (HGBT, $q=10$)}
\label{tab:exp:DSB:HGBT:q10}
\centering\small
\resizebox{\textwidth}{!}{%
\begin{tabular}{lrrrrrr}
\hline
Метод & ROC AUC & PR AUC & $\Delta$ROC AUC & $\Delta$PR AUC & $t_{\text{ред}}$, c & $t_{\text{мод}}$, c\\
\hline
Ensembleunar & 0.959129 & 0.906669 & 0.000708 & 0.000754 & 1466.362532 & 0.340065\\
PLS          & 0.954915 & 0.902382 & -0.003506 & -0.003532 & 0.012408 & 0.319568\\
HSIC\_Lasso  & 0.944929 & 0.878895 & -0.013491 & -0.027020 & 0.699972 & 0.317398\\
PCA          & 0.876723 & 0.688229 & -0.081698 & -0.217685 & 0.002740 & 0.337652\\
UMAP\_sup    & 0.775408 & 0.567177 & -0.183012 & -0.338737 & 20.492583 & 0.341360\\
\hline
\end{tabular}}
\end{table}

\begin{table}[H]
\caption{Результаты на DS-B (HGBT, $q=15$)}
\label{tab:exp:DSB:HGBT:q15}
\centering\small
\resizebox{\textwidth}{!}{%
\begin{tabular}{lrrrrrr}
\hline
Метод & ROC AUC & PR AUC & $\Delta$ROC AUC & $\Delta$PR AUC & $t_{\text{ред}}$, c & $t_{\text{мод}}$, c\\
\hline
Ensembleunar & 0.955862 & 0.901995 & -0.002558 & -0.003919 & 1469.874476 & 0.344023\\
PLS          & 0.952853 & 0.899741 & -0.005567 & -0.006173 & 0.016798 & 0.335751\\
HSIC\_Lasso  & 0.942160 & 0.876210 & -0.016260 & -0.029704 & 0.936248 & 0.365223\\
PCA          & 0.889813 & 0.716948 & -0.068608 & -0.188967 & 0.003392 & 0.331456\\
UMAP\_sup    & 0.782580 & 0.573703 & -0.175840 & -0.332212 & 23.408102 & 0.360922\\
\hline
\end{tabular}}
\end{table}

% ---- DS-B, MLP ----
\begin{table}[H]
\caption{Результаты на DS-B (MLP, $q=5$)}
\label{tab:exp:DSB:MLP:q5}
\centering\small
\resizebox{\textwidth}{!}{%
\begin{tabular}{lrrrrrr}
\hline
Метод & ROC AUC & PR AUC & $\Delta$ROC AUC & $\Delta$PR AUC & $t_{\text{ред}}$, c & $t_{\text{мод}}$, c\\
\hline
Ensembleunar & 0.960336 & 0.914073 & 0.028133 & 0.049714 & 1462.708774 & 1.023030\\
HSIC\_Lasso  & 0.953842 & 0.906015 & 0.021639 & 0.041656 & 0.605554 & 0.630133\\
PLS          & 0.939943 & 0.872374 & 0.007740 & 0.008015 & 0.007177 & 0.712632\\
PCA          & 0.687974 & 0.448327 & -0.244230 & -0.416032 & 0.002488 & 0.397214\\
UMAP\_sup    & 0.549548 & 0.289236 & -0.382655 & -0.575123 & 16.514097 & 0.179547\\
\hline
\end{tabular}}
\end{table}

\begin{table}[H]
\caption{Результаты на DS-B (MLP, $q=10$)}
\label{tab:exp:DSB:MLP:q10}
\centering\small
\resizebox{\textwidth}{!}{%
\begin{tabular}{lrrrrrr}
\hline
Метод & ROC AUC & PR AUC & $\Delta$ROC AUC & $\Delta$PR AUC & $t_{\text{ред}}$, c & $t_{\text{мод}}$, c\\
\hline
Ensembleunar & 0.964891 & 0.921808 & 0.032688 & 0.057449 & 1466.362532 & 1.074756\\
PLS          & 0.953977 & 0.900823 & 0.021773 & 0.036464 & 0.012408 & 0.645237\\
HSIC\_Lasso  & 0.939133 & 0.876846 & 0.006929 & 0.012487 & 0.699972 & 0.762155\\
PCA          & 0.868926 & 0.681057 & -0.063277 & -0.183303 & 0.002740 & 0.695374\\
UMAP\_sup    & 0.484301 & 0.249009 & -0.447902 & -0.615350 & 20.492583 & 0.178839\\
\hline
\end{tabular}}
\end{table}

\begin{table}[H]
\caption{Результаты на DS-B (MLP, $q=15$)}
\label{tab:exp:DSB:MLP:q15}
\centering\small
\resizebox{\textwidth}{!}{%
\begin{tabular}{lrrrrrr}
\hline
Метод & ROC AUC & PR AUC & $\Delta$ROC AUC & $\Delta$PR AUC & $t_{\text{ред}}$, c & $t_{\text{мод}}$, c\\
\hline
Ensembleunar & 0.944871 & 0.888826 & 0.012668 & 0.024467 & 1469.874476 & 0.752957\\
PLS          & 0.944392 & 0.884113 & 0.012189 & 0.019754 & 0.016798 & 0.592478\\
HSIC\_Lasso  & 0.935508 & 0.870047 & 0.003305 & 0.005688 & 0.936248 & 1.190793\\
PCA          & 0.893645 & 0.735038 & -0.038558 & -0.129321 & 0.003392 & 0.747590\\
UMAP\_sup    & 0.548465 & 0.290475 & -0.383738 & -0.573884 & 23.408102 & 0.178606\\
\hline
\end{tabular}}
\end{table}

На DS-B метод Ensembleunar демонстрирует наилучшее качество среди сравниваемых подходов
для обеих последующих моделей (HGBT и MLP) при $q\in\{10,15\}$, а для MLP также при $q=5$.
При $q=5$ и HGBT метод уступает HSIC\_Lasso, однако разница по ROC AUC порядка $0.01$ остаётся в пределах сопоставимого качества.
Существенным практическим ограничением является высокая вычислительная стоимость редуктора по времени.

% ----------------------------------------------------------------------
\subsection{Реальный набор данных DS-A-real}

В качестве первого реального набора данных используется набор данных Rice (Cammeo and Osmancik),
представляющий собой задачу бинарной классификации двух сортов риса на основе морфологических
характеристик отдельных зёрен. Каждый объект описывается $d=7$ компактными геометрическими
признаками, полученными в результате обработки изображений зёрен, включая характеристики формы,
размеров и контура. Эти признаки хорошо интерпретируются с точки зрения предметной области и при
этом содержат достаточную различительную информацию между сортами, что делает данный набор данных
удобным тестовым примером для методов отбора признаков и понижения размерности в условиях
относительно малой размерности и умеренной коррелированности координат. Описание набора данных
и исходная постановка задачи приведены в работе \cite{uci_rice_2019}.
В экспериментах для набора данных DS-A-real рассматривались размеры редукции $q\in\{2,3,5\}$.

% ---- DS-A-real, HGBT ----
\begin{table}[H]
\caption{Результаты на DS-A-real (HGBT, $q=2$)}
\label{tab:exp:DSAreal:HGBT:q2}
\centering\small
\resizebox{\textwidth}{!}{%
\begin{tabular}{lrrrrrr}
\hline
Метод & ROC AUC & PR AUC & $\Delta$ROC AUC & $\Delta$PR AUC & $t_{\text{ред}}$, c & $t_{\text{мод}}$, c\\
\hline
HSIC\_Lasso & 0.973615 & 0.979031 & 0.000137 & -0.000615 & 0.323406 & 0.200964\\
MLP\_unar   & 0.968427 & 0.974582 & -0.005051 & -0.005064 & 312.459187 & 0.233460\\
PCA         & 0.969744 & 0.976321 & -0.003733 & -0.003325 & 0.001330 & 0.202325\\
PLS         & 0.973056 & 0.979334 & -0.000421 & -0.000312 & 0.002072 & 0.204862\\
UMAP\_sup   & 0.969421 & 0.975281 & -0.004057 & -0.004364 & 17.879416 & 0.203783\\
\hline
\end{tabular}}
\end{table}

\begin{table}[H]
\caption{Результаты на DS-A-real (HGBT, $q=3$)}
\label{tab:exp:DSAreal:HGBT:q3}
\centering\small
\resizebox{\textwidth}{!}{%
\begin{tabular}{lrrrrrr}
\hline
Метод & ROC AUC & PR AUC & $\Delta$ROC AUC & $\Delta$PR AUC & $t_{\text{ред}}$, c & $t_{\text{мод}}$, c\\
\hline
HSIC\_Lasso & 0.973096 & 0.979010 & -0.000381 & -0.000636 & 0.326939 & 0.211409\\
MLP\_unar   & 0.972266 & 0.979232 & -0.001212 & -0.000413 & 342.233841 & 0.296109\\
PCA         & 0.971694 & 0.977297 & -0.001784 & -0.002349 & 0.001181 & 0.216057\\
PLS         & 0.973231 & 0.979678 & -0.000246 & 0.000032 & 0.002241 & 0.211898\\
UMAP\_sup   & 0.970269 & 0.977100 & -0.003208 & -0.002546 & 17.608202 & 0.213535\\
\hline
\end{tabular}}
\end{table}

\begin{table}[H]
\caption{Результаты на DS-A-real (HGBT, $q=5$)}
\label{tab:exp:DSAreal:HGBT:q5}
\centering\small
\resizebox{\textwidth}{!}{%
\begin{tabular}{lrrrrrr}
\hline
Метод & ROC AUC & PR AUC & $\Delta$ROC AUC & $\Delta$PR AUC & $t_{\text{ред}}$, c & $t_{\text{мод}}$, c\\
\hline
HSIC\_Lasso & 0.973252 & 0.979620 & -0.000226 & -0.000026 & 0.332584 & 0.223249\\
MLP\_unar   & 0.973243 & 0.979753 & -0.000235 & 0.000107 & 159.219546 & 0.223269\\
PCA         & 0.972277 & 0.977955 & -0.001200 & -0.001691 & 0.001325 & 0.239144\\
PLS         & 0.974208 & 0.980271 & 0.000730 & 0.000625 & 0.002769 & 0.228442\\
UMAP\_sup   & 0.970365 & 0.976384 & -0.003113 & -0.003261 & 18.400578 & 0.237983\\
\hline
\end{tabular}}
\end{table}

% ---- DS-A-real, MLP ----
\begin{table}[H]
\caption{Результаты на DS-A-real (MLP, $q=2$)}
\label{tab:exp:DSAreal:MLP:q2}
\centering\small
\resizebox{\textwidth}{!}{%
\begin{tabular}{lrrrrrr}
\hline
Метод & ROC AUC & PR AUC & $\Delta$ROC AUC & $\Delta$PR AUC & $t_{\text{ред}}$, c & $t_{\text{мод}}$, c\\
\hline
HSIC\_Lasso & 0.977700 & 0.982612 & 0.000367 & 0.000213 & 0.323406 & 0.331590\\
MLP\_unar   & 0.973095 & 0.977879 & -0.004238 & -0.004520 & 312.459187 & 0.301596\\
PCA         & 0.976083 & 0.981470 & -0.001250 & -0.000929 & 0.001330 & 0.241261\\
PLS         & 0.977713 & 0.982632 & 0.000380 & 0.000233 & 0.002072 & 0.224398\\
UMAP\_sup   & 0.965429 & 0.969389 & -0.011904 & -0.013010 & 17.879416 & 0.459461\\
\hline
\end{tabular}}
\end{table}

\begin{table}[H]
\caption{Результаты на DS-A-real (MLP, $q=3$)}
\label{tab:exp:DSAreal:MLP:q3}
\centering\small
\resizebox{\textwidth}{!}{%
\begin{tabular}{lrrrrrr}
\hline
Метод & ROC AUC & PR AUC & $\Delta$ROC AUC & $\Delta$PR AUC & $t_{\text{ред}}$, c & $t_{\text{мод}}$, c\\
\hline
HSIC\_Lasso & 0.977711 & 0.983051 & 0.000379 & 0.000652 & 0.326939 & 0.280160\\
MLP\_unar   & 0.977507 & 0.983069 & 0.000174 & 0.000670 & 342.233841 & 0.347544\\
PCA         & 0.977507 & 0.982776 & 0.000174 & 0.000378 & 0.001181 & 0.186016\\
PLS         & 0.977316 & 0.982379 & -0.000017 & -0.000020 & 0.002241 & 0.200615\\
UMAP\_sup   & 0.968809 & 0.973760 & -0.008524 & -0.008639 & 17.608202 & 0.322675\\
\hline
\end{tabular}}
\end{table}

\begin{table}[H]
\caption{Результаты на DS-A-real (MLP, $q=5$)}
\label{tab:exp:DSAreal:MLP:q5}
\centering\small
\resizebox{\textwidth}{!}{%
\begin{tabular}{lrrrrrr}
\hline
Метод & ROC AUC & PR AUC & $\Delta$ROC AUC & $\Delta$PR AUC & $t_{\text{ред}}$, c & $t_{\text{мод}}$, c\\
\hline
HSIC\_Lasso & 0.976299 & 0.981623 & -0.001034 & -0.000776 & 0.332584 & 0.205971\\
MLP\_unar   & 0.976279 & 0.981302 & -0.001054 & -0.001097 & 159.219546 & 0.247190\\
PCA         & 0.975135 & 0.980608 & -0.002198 & -0.001791 & 0.001325 & 0.277232\\
PLS         & 0.978093 & 0.983197 & 0.000760 & 0.000798 & 0.002769 & 0.298271\\
UMAP\_sup   & 0.971236 & 0.976596 & -0.006097 & -0.005803 & 18.400578 & 0.239608\\
\hline
\end{tabular}}
\end{table}

На DS-A-real метод MLP\_unar находится на уровне HSIC\_Lasso и близок к линейным редукторам (PCA/PLS),
а отставание от лучшего результата, как правило, составляет тысячные.
При $q=2$ наблюдается более заметный разрыв по сравнению с лидерами, однако он остаётся в пределах $\approx 0.005$ по ROC AUC,
то есть качество остаётся сопоставимым; различия могут зависеть от режима валидации и настройки гиперпараметров.

% ----------------------------------------------------------------------
\subsection{Реальный набор данных DS-B-real}

В качестве второго реального набора данных используется набор данных 
Phishing Websites, который включает $d=30$ табличных признаков, 
характеризующих URL и свойства веб-страниц, и предназначен для задачи 
бинарной классификации фишинговых сайтов. Признаки в наборе данных дискретны и 
отражают различные аспекты, такие как длина URL, использование поддоменов, 
наличие специальных символов и другие эвристики, позволяющие отличить 
фишинговые сайты от легитимных. Этот набор данных моделирует практическую задачу 
в области кибербезопасности, где разделимость классов достигается за счет 
комбинации различных признаков. Описание набора данных приведено в работе 
\cite{uci_phishing_2012}. В экспериментах для набора данных 
DS-B-real использовались размеры редукции $q\in\{5,10,15\}$.

% ---- DS-B-real, HGBT ----
\begin{table}[H]
\caption{Результаты на DS-B-real (HGBT, $q=5$)}
\label{tab:exp:DSBreal:HGBT:q5}
\centering\small
\resizebox{\textwidth}{!}{%
\begin{tabular}{lrrrrrr}
\hline
Метод & ROC AUC & PR AUC & $\Delta$ROC AUC & $\Delta$PR AUC & $t_{\text{ред}}$, c & $t_{\text{мод}}$, c\\
\hline
PLS         & 0.989753 & 0.991517 & -0.006066 & -0.004977 & 0.012622 & 0.319731\\
PCA         & 0.986859 & 0.988822 & -0.008960 & -0.007672 & 0.003560 & 0.316847\\
HSIC\_Lasso & 0.981596 & 0.982182 & -0.014223 & -0.014313 & 1.201798 & 0.320936\\
Ensembleunar& 0.970599 & 0.968780 & -0.025220 & -0.027714 & 434.225702 & 0.294692\\
UMAP\_sup   & 0.967280 & 0.972473 & -0.028539 & -0.024021 & 21.764809 & 0.330768\\
\hline
\end{tabular}}
\end{table}

\begin{table}[H]
\caption{Результаты на DS-B-real (HGBT, $q=10$)}
\label{tab:exp:DSBreal:HGBT:q10}
\centering\small
\resizebox{\textwidth}{!}{%
\begin{tabular}{lrrrrrr}
\hline
Метод & ROC AUC & PR AUC & $\Delta$ROC AUC & $\Delta$PR AUC & $t_{\text{ред}}$, c & $t_{\text{мод}}$, c\\
\hline
PCA         & 0.991934 & 0.993346 & -0.003885 & -0.003148 & 0.012073 & 0.345465\\
PLS         & 0.991421 & 0.992924 & -0.004398 & -0.003570 & 0.019920 & 0.336661\\
HSIC\_Lasso & 0.990009 & 0.991589 & -0.005810 & -0.004906 & 1.393820 & 0.335443\\
Ensembleunar& 0.981221 & 0.983066 & -0.014597 & -0.013429 & 434.457498 & 0.319707\\
UMAP\_sup   & 0.971977 & 0.976625 & -0.023842 & -0.019869 & 23.361637 & 0.351070\\
\hline
\end{tabular}}
\end{table}

\begin{table}[H]
\caption{Результаты на DS-B-real (HGBT, $q=15$)}
\label{tab:exp:DSBreal:HGBT:q15}
\centering\small
\resizebox{\textwidth}{!}{%
\begin{tabular}{lrrrrrr}
\hline
Метод & ROC AUC & PR AUC & $\Delta$ROC AUC & $\Delta$PR AUC & $t_{\text{ред}}$, c & $t_{\text{мод}}$, c\\
\hline
PLS         & 0.993227 & 0.994478 & -0.002592 & -0.002016 & 0.026245 & 0.369544\\
PCA         & 0.993068 & 0.994353 & -0.002751 & -0.002141 & 0.012378 & 0.378682\\
HSIC\_Lasso & 0.990512 & 0.992107 & -0.005306 & -0.004388 & 1.392486 & 0.334822\\
Ensembleunar& 0.987567 & 0.989203 & -0.008252 & -0.007291 & 438.116849 & 0.336248\\
UMAP\_sup   & 0.972000 & 0.976483 & -0.023818 & -0.020011 & 31.216013 & 0.359395\\
\hline
\end{tabular}}
\end{table}

% ---- DS-B-real, MLP ----
\begin{table}[H]
\caption{Результаты на DS-B-real (MLP, $q=5$)}
\label{tab:exp:DSBreal:MLP:q5}
\centering\small
\resizebox{\textwidth}{!}{%
\begin{tabular}{lrrrrrr}
\hline
Метод & ROC AUC & PR AUC & $\Delta$ROC AUC & $\Delta$PR AUC & $t_{\text{ред}}$, c & $t_{\text{мод}}$, c\\
\hline
PLS         & 0.978977 & 0.982998 & -0.015351 & -0.012374 & 0.012622 & 0.925834\\
HSIC\_Lasso & 0.977065 & 0.978955 & -0.017262 & -0.016418 & 1.201798 & 1.208000\\
PCA         & 0.970663 & 0.975464 & -0.023664 & -0.019909 & 0.003560 & 0.612417\\
Ensembleunar& 0.964499 & 0.963916 & -0.029828 & -0.031457 & 434.225702 & 0.618334\\
UMAP\_sup   & 0.913640 & 0.930414 & -0.080688 & -0.064959 & 21.764809 & 2.222381\\
\hline
\end{tabular}}
\end{table}

\begin{table}[H]
\caption{Результаты на DS-B-real (MLP, $q=10$)}
\label{tab:exp:DSBreal:MLP:q10}
\centering\small
\resizebox{\textwidth}{!}{%
\begin{tabular}{lrrrrrr}
\hline
Метод & ROC AUC & PR AUC & $\Delta$ROC AUC & $\Delta$PR AUC & $t_{\text{ред}}$, c & $t_{\text{мод}}$, c\\
\hline
HSIC\_Lasso & 0.984309 & 0.987019 & -0.010019 & -0.008353 & 1.393820 & 1.063420\\
PCA         & 0.984345 & 0.987509 & -0.009982 & -0.007864 & 0.012073 & 1.321070\\
PLS         & 0.983363 & 0.986768 & -0.010964 & -0.008605 & 0.019920 & 1.041679\\
Ensembleunar& 0.975694 & 0.978179 & -0.018633 & -0.017194 & 434.457498 & 1.045043\\
UMAP\_sup   & 0.917282 & 0.933396 & -0.077045 & -0.061977 & 23.361637 & 1.782164\\
\hline
\end{tabular}}
\end{table}

\begin{table}[H]
\caption{Результаты на DS-B-real (MLP, $q=15$)}
\label{tab:exp:DSBreal:MLP:q15}
\centering\small
\resizebox{\textwidth}{!}{%
\begin{tabular}{lrrrrrr}
\hline
Метод & ROC AUC & PR AUC & $\Delta$ROC AUC & $\Delta$PR AUC & $t_{\text{ред}}$, c & $t_{\text{мод}}$, c\\
\hline
PCA         & 0.988880 & 0.991176 & -0.005447 & -0.004197 & 0.012378 & 1.876479\\
PLS         & 0.988173 & 0.990636 & -0.006155 & -0.004736 & 0.026245 & 1.738431\\
HSIC\_Lasso & 0.986386 & 0.988938 & -0.007942 & -0.006435 & 1.392486 & 1.553514\\
Ensembleunar& 0.982562 & 0.984832 & -0.011765 & -0.010540 & 438.116849 & 0.996738\\
UMAP\_sup   & 0.900993 & 0.920600 & -0.093335 & -0.074773 & 31.216013 & 1.602632\\
\hline
\end{tabular}}
\end{table}

На DS-B-real метод Ensembleunar уступает HSIC\_Lasso,
однако разница по ROC AUC во всех рассмотренных режимах не превышает $\approx 0.02$,
то есть результаты остаются сопоставимыми по качеству и могут зависеть от режима настройки
гиперпараметров и статистических свойств процедуры оценки информативности.
Однако при этом вычислительная стоимость редуктора существенно выше.

\section*{Заключение}
\addcontentsline{toc}{section}{Заключение}

В работе предложены и исследованы два новых метода оценки информативности признаков табличных данных
в задаче бинарной классификации, основанные на идее унарной классификации с добавлением фонового класса,
равномерно распределённого на компакте. Такая постановка фиксирует корректное поведение модели вне носителя
данных и формирует естественную зону отказа, что делает последующие оценки информативности более устойчивыми.

Первый метод представляет собой адаптацию схемы последовательного исключения признаков с переобучением моделей 
к унарной постановке и использует унарные метрики разделимости, учитывающие области
однозначного решения, наложения и отказа. Второй метод сочетает ансамбли случайных подпространств 
с унарными классификаторами и обновлением распределения важностей 
координат на основе аддитивных объяснений Шепли, что позволяет частично преодолевать проклятие размерности
за счёт обучения на подпространствах малой размерности.

Экспериментальное исследование проведено на четырёх наборах данных: двух синтетических и двух реальных. Предложенные методы сравнивались с базовыми подходами понижения
размерности и отбора признаков, включая метод главных компонент, метод частичных наименьших квадратов,
проекцию многообразия в постановке с учителем, а также отбор признаков по критерию независимости
Гильберта-Шмидта с разреживающей регуляризацией. Качество оценивалось по влиянию редукции на последующие модели
классификации в терминах ROC AUC и PR AUC.

Полученные результаты показывают конкурентоспособность предложенных подходов. 
В дальнейших работах планируется улучшение результатов за счёт систематического подбора гиперпараметров унарной части,
а также за счёт проведения более широких экспериментов на дополнительных реальных наборах данных и в различных режимах дисбаланса классов.
Отдельным направлением является усиление теоретического обоснования предложенных методов.
Кроме того, предполагается модифицировать критерии выбора моделей и порогов в унарной постановке так, чтобы перенести предложенные идеи
на многоклассовый случай.

Таким образом, уже на текущем этапе предложенные методы формируют работоспособный и интерпретируемый инструментарий оценки информативности
признаков в табличной классификации на базе унарного подхода с зоной отказа, демонстрируя качество на уровне сильных базовых методов
и задавая направления для дальнейшего повышения точности, устойчивости и практической эффективности.

\addcontentsline{toc}{section}{Список использованных источников}
\begin{thebibliography}{99}
    \bibitem{unar} Lukyanov K. S. et al. Extrapolation of the Bayesian classifier with an unknown support of the two-class mixture distribution //Uspekhi Matematicheskikh Nauk. – 2024. – Т. 79. – №. 6. – С. 57-82.
    \bibitem{Pearson1901} Pearson, K. On lines and planes of closest fit to systems of points in space. Philosophical Magazine, Series 6, 2(11), 559--572 (1901).
    \bibitem{Hotelling1933} Hotelling, H. Analysis of a complex of statistical variables into principal components. Journal of Educational Psychology, 24, 417--441, 498--520 (1933).
    \bibitem{WoldRuhe1984} Wold, S., Ruhe, A., Wold, H., Dunn, W. J. The collinearity problem in linear regression. The partial least squares approach to generalized inverses. SIAM Journal on Scientific and Statistical Computing, 5(3), 735--743 (1984).
    \bibitem{McInnes2018arXiv} McInnes, L., Healy, J., Melville, J. UMAP: Uniform Manifold Approximation and Projection for Dimension Reduction. arXiv:1802.03426 (2018).
    \bibitem{McInnes2018JOSS} McInnes, L., Healy, J., Saul, N., Grossberger, L. UMAP: Uniform Manifold Approximation and Projection. Journal of Open Source Software, 3(29), 861 (2018).
    \bibitem{Gretton2005} Gretton, A., Bousquet, O., Smola, A., Sch{\"o}lkopf, B. Measuring Statistical Dependence with Hilbert-Schmidt Norms. In: Algorithmic Learning Theory (ALT), 63--77 (2005).
    \bibitem{Yamada2014} Yamada, M., Jitkrittum, W., Sigal, L., Xing, E. P., Sugiyama, M. High-Dimensional Feature Selection by Feature-Wise Kernelized Lasso. Neural Computation, 26(1), 185--207 (2014).
    \bibitem{eliseev2025} Н.\,А.~Елисеев, А.\,И.~Перминов, Д.\,Ю.~Турдаков, Сходимость многослойного персептрона к гистограммной байесовской регрессии, УМН, 80:6(486) (2025), 45--72.
    \bibitem{liu2013lofo} Liu~J., Danait~N., Hu~S., Sengupta~S. A leave-one-feature-out wrapper method for feature selection in data classification // Proceedings of the 2013 6th International Conference on Biomedical Engineering and Informatics (BMEI). 2013. DOI: 10.1109/BMEI.2013.6747021.
    \bibitem{tian2021rase} Tian~Y., Feng~Y. RaSE: Random subspace ensemble classification // Journal of Machine Learning Research. --- 2021. --- Т.~22. --- №~45. --- С.~1--93.
    \bibitem{lundberg2017shap}Lundberg~S.~M., Lee~S.-I. A Unified Approach to Interpreting Model Predictions // Advances in Neural Information Processing Systems. --- 2017. --- Vol.~30. --- P.~4765--4774.
    \bibitem{uci_rice_2019} Rice (Cammeo and Osmancik) [Dataset]. UCI Machine Learning Repository, 2019. DOI: 10.24432/C5MW4Z.
    \bibitem{uci_phishing_2012} Mohammad~R., McCluskey~L. Phishing Websites [Dataset]. UCI Machine Learning Repository, 2012. DOI: 10.24432/C51W2X.
\end{thebibliography}
\section*{ПРИЛОЖЕНИЯ}
\addcontentsline{toc}{section}{ПРИЛОЖЕНИЯ}

Ниже приведен код всех экспериментов и методов, описанных в работе.

\subsection*{ПРИЛОЖЕНИЕ А}
\addcontentsline{toc}{subsection}{ПРИЛОЖЕНИЕ А}

Ниже приведен код, относящийся к LOFO методу и частично используемый для RaSE метода.

\begin{lstlisting}[
    language=Python,
    caption={Метод LOFO с использованием унарной классификации},
    label={lst:process_datasets},
    captionpos=b
]
import numpy as np
import pandas as pd
from time import perf_counter
from typing import Any, Dict, List

from sklearn.datasets import make_classification
from sklearn.model_selection import StratifiedKFold
from sklearn.preprocessing import StandardScaler
from sklearn.metrics import roc_auc_score, average_precision_score, precision_recall_fscore_support
from sklearn.neural_network import MLPClassifier
from sklearn.ensemble import HistGradientBoostingClassifier

from sklearn.decomposition import PCA
from sklearn.cross_decomposition import PLSRegression
from umap import UMAP

from pyHSICLasso import HSICLasso

import torch
import torch.nn as nn
import torch.nn.functional as F
from torch.utils.data import DataLoader, TensorDataset, random_split, Subset, Dataset, ConcatDataset
from sklearn.model_selection import train_test_split
from typing import Tuple, Optional, Iterable
import random
import shap
import copy
import pickle  
from sklearn.model_selection import StratifiedShuffleSplit
from ucimlrepo import fetch_ucirepo

def compute_feature_bounds(dataset: Dataset, pad: float = 0.25):
    if isinstance(dataset, torch.utils.data.TensorDataset):
        X = dataset.tensors[0].to(dtype=torch.float32)
        fmin = X.amin(dim=0)
        fmax = X.amax(dim=0)
        span = fmax - fmin
        return fmin - pad * span, fmax + pad * span

    if isinstance(dataset, torch.utils.data.Subset) and isinstance(dataset.dataset, torch.utils.data.TensorDataset):
        Xfull = dataset.dataset.tensors[0]
        idx = torch.as_tensor(dataset.indices, device=Xfull.device, dtype=torch.long)
        X = Xfull.index_select(0, idx).to(dtype=torch.float32)
        fmin = X.amin(dim=0)
        fmax = X.amax(dim=0)
        span = fmax - fmin
        return fmin - pad * span, fmax + pad * span

    xs = [dataset[i][0] for i in range(len(dataset))]
    if len(xs) == 0:
        raise ValueError("compute_feature_bounds: empty dataset")
    device = xs[0].device if torch.is_tensor(xs[0]) else torch.device("cpu")
    xs = [x.to(device) if torch.is_tensor(x) else torch.as_tensor(x, device=device) for x in xs]
    X = torch.stack(xs, dim=0).to(dtype=torch.float32)
    fmin = X.min(dim=0).values
    fmax = X.max(dim=0).values
    span = fmax - fmin
    return fmin - pad * span, fmax + pad * span

class UniformNoiseDataset(Dataset):
    def __init__(
        self,
        feature_min: torch.Tensor,
        feature_max: torch.Tensor,
        n: int,
        label: int = 0,
        dtype: torch.dtype | None = None
    ):
        assert feature_min.shape == feature_max.shape

        self.dtype = dtype if dtype is not None else torch.float32

        feature_min = feature_min.to(dtype=self.dtype)
        feature_max = feature_max.to(device=feature_min.device, dtype=self.dtype)

        self.feature_min = feature_min
        self.feature_max = feature_max
        self.n = int(n)
        self.label = int(label)

    def __len__(self) -> int:
        return self.n

    def __getitem__(self, idx: int):
        u = torch.rand_like(self.feature_min, dtype=self.dtype)
        x = self.feature_min + u * (self.feature_max - self.feature_min)
        y = torch.tensor(self.label, dtype=torch.long, device=x.device)
        return x, y

class LabelMappedSubset(Dataset):
    def __init__(self, subset: Subset, label_map: dict[int, int]):
        self.subset = subset
        self.label_map = label_map

    def __getitem__(self, idx):
        x, y = self.subset[idx]

        if torch.is_tensor(y):
            y_int = int(y.detach().cpu().item())
        else:
            y_int = int(y)

        y_mapped = self.label_map.get(y_int, y_int)
        return x, torch.tensor(y_mapped, dtype=y.dtype, device=y.device)

    def __len__(self):
        return len(self.subset)


class Abs(nn.Module):
    def forward(self, x: torch.Tensor) -> torch.Tensor:
        return torch.abs(x)

class MLP(nn.Module):
    def __init__(self, d_in: int, d_hidden: int = 10, n_hidden_layers: int = 2):

        super().__init__()

        layers = []

        layers.append(nn.Linear(d_in, d_hidden))
        layers.append(Abs())


        for _ in range(n_hidden_layers - 1):
            layers.append(nn.Linear(d_hidden, d_hidden))
            layers.append(Abs())

        layers.append(nn.Linear(d_hidden, 1))

        self.net = nn.Sequential(*layers)

    def forward(self, x: torch.Tensor) -> torch.Tensor:
        return self.net(x)

def train_one_model(
    model: nn.Module,
    dataset_pos_as_1: Dataset,
    feature_min: torch.Tensor,
    feature_max: torch.Tensor,
    num_epochs: int = 20,
    batch_size: int = 128,
    lr: float = 3e-3,
    weight_decay: float = 1e-2,
    seed: int = 0,
):
    """
    Обучает модель различать:
      - реальные (label=1)
      - шумовые (label=0)
    с MSELoss.
    """
    torch.manual_seed(seed)
    if torch.cuda.is_available():
        torch.cuda.manual_seed_all(seed)

    device = next(model.parameters()).device
    criterion = nn.MSELoss()
    opt = torch.optim.AdamW(model.parameters(), lr=lr, weight_decay=weight_decay)

    if isinstance(dataset_pos_as_1, TensorDataset):
        X_real = dataset_pos_as_1.tensors[0].to(device=device, dtype=torch.float32)
        n = int(X_real.shape[0])
        if n == 0:
            raise ValueError("train_one_model: empty TensorDataset")

        fmin = feature_min.to(device=device, dtype=torch.float32)
        fmax = feature_max.to(device=device, dtype=torch.float32)
        span = (fmax - fmin)
        d = int(X_real.shape[1])

        for _ in range(num_epochs):
            model.train()
            perm = torch.randperm(n, device=device)

            for start in range(0, n, batch_size):
                idx = perm[start:start + batch_size]
                xb_real = X_real.index_select(0, idx)


                u = torch.rand(xb_real.shape[0], d, device=device, dtype=torch.float32)
                xb_noise = fmin.unsqueeze(0) + u * span.unsqueeze(0) 

                xb = torch.cat([xb_real, xb_noise], dim=0) 
                yb = torch.cat(
                    [
                        torch.ones(xb_real.shape[0], device=device, dtype=torch.float32),
                        torch.zeros(xb_noise.shape[0], device=device, dtype=torch.float32),
                    ],
                    dim=0,
                ).unsqueeze(1) 

                pred = model(xb)
                loss = criterion(pred, yb)

                opt.zero_grad(set_to_none=True)
                loss.backward()
                opt.step()

        return model


    n_real = len(dataset_pos_as_1)
    if n_real <= 0:
        raise ValueError("train_one_model: empty dataset_pos_as_1")

    n_noise = max(1, n_real)

    x0, _ = dataset_pos_as_1[0]
    data_device = x0.device if torch.is_tensor(x0) else device

    fmin = feature_min.detach().to(device=data_device, dtype=torch.float32)
    fmax = feature_max.detach().to(device=data_device, dtype=torch.float32)

    noise_ds = UniformNoiseDataset(fmin, fmax, n_noise, label=0, dtype=torch.float32)
    train_ds = ConcatDataset([dataset_pos_as_1, noise_ds])
    loader = DataLoader(train_ds, batch_size=batch_size, shuffle=True, drop_last=False)

    for _ in range(num_epochs):
        model.train()
        for xb, yb in loader:
            xb = xb.to(device=device, dtype=torch.float32)
            yb = yb.to(device=device).unsqueeze(-1).to(dtype=torch.float32)

            pred = model(xb)
            loss = criterion(pred, yb)

            opt.zero_grad(set_to_none=True)
            loss.backward()
            opt.step()

    return model

def model_train(
    d_hidden: int,
    n_hidden_layers: int,
    train_data: Dataset,
    num_epochs: int,
    batch_size: int = 128,
    seed: int = 0,
    output: bool = False,
    device: str | torch.device | None = None,
    feature_min: torch.Tensor | None = None,
    feature_max: torch.Tensor | None = None,
):
    torch.manual_seed(seed)
    if torch.cuda.is_available():
        torch.cuda.manual_seed_all(seed)

    if isinstance(train_data, TensorDataset):
        X_all, y_all = train_data.tensors


        data_device = X_all.device
        dev = data_device if device is None else torch.device(device)

        X_all = X_all.to(device=dev, dtype=torch.float32)
        y_all = y_all.to(device=dev, dtype=torch.long)

        d_in = int(X_all.shape[1])

        if feature_min is None or feature_max is None:
            fmin, fmax = compute_feature_bounds(TensorDataset(X_all, y_all))
        else:
            fmin, fmax = feature_min, feature_max
        fmin = fmin.to(device=dev, dtype=torch.float32)
        fmax = fmax.to(device=dev, dtype=torch.float32)

        pos_mask = (y_all == 1)
        neg_mask = (y_all == -1)
        X_pos = X_all[pos_mask]
        X_neg = X_all[neg_mask]

        if X_pos.numel() == 0 or X_neg.numel() == 0:
            raise ValueError("model_train: empty class after split (pos or neg)")

        model_pos = MLP(d_in=d_in, d_hidden=d_hidden, n_hidden_layers=n_hidden_layers).to(dev)
        model_neg = MLP(d_in=d_in, d_hidden=d_hidden, n_hidden_layers=n_hidden_layers).to(dev)

        ds_pos = TensorDataset(X_pos, torch.ones(X_pos.size(0), device=dev, dtype=torch.long))
        ds_neg = TensorDataset(X_neg, torch.ones(X_neg.size(0), device=dev, dtype=torch.long))

        model_pos = train_one_model(
            model_pos, ds_pos, fmin, fmax,
            num_epochs=num_epochs, batch_size=batch_size, seed=seed
        )
        model_neg = train_one_model(
            model_neg, ds_neg, fmin, fmax,
            num_epochs=num_epochs, batch_size=batch_size, seed=seed + 1
        )
        return model_pos, model_neg

    ys: list[int] = []
    for i in range(len(train_data)):
        y = train_data[i][1]
        if torch.is_tensor(y):
            ys.append(int(y.detach().cpu().item()))
        else:
            ys.append(int(y))
    labels = torch.tensor(ys, dtype=torch.long)

    pos_idx = (labels == 1).nonzero(as_tuple=True)[0].tolist()
    neg_idx = (labels == -1).nonzero(as_tuple=True)[0].tolist()

    train_pos = Subset(train_data, pos_idx)
    train_neg = Subset(train_data, neg_idx)
    train_neg_as_pos = LabelMappedSubset(train_neg, {-1: 1})

    x0, _ = train_data[0]
    d_in = x0.shape[-1]
    data_device = x0.device if torch.is_tensor(x0) else torch.device("cpu")

    dev = data_device if device is None else torch.device(device)

    fmin, fmax = compute_feature_bounds(train_data)
    fmin = fmin.to(dev)
    fmax = fmax.to(dev)

    model_pos = MLP(d_in=d_in, d_hidden=d_hidden, n_hidden_layers=n_hidden_layers).to(dev)
    model_neg = MLP(d_in=d_in, d_hidden=d_hidden, n_hidden_layers=n_hidden_layers).to(dev)

    model_pos = train_one_model(
        model_pos, train_pos, fmin, fmax,
        num_epochs=num_epochs, batch_size=batch_size, seed=seed
    )
    model_neg = train_one_model(
        model_neg, train_neg_as_pos, fmin, fmax,
        num_epochs=num_epochs, batch_size=batch_size, seed=seed + 1
    )
    return model_pos, model_neg

def metrics_from_sheet(X, y, model_pos, model_neg, beta_pos=0.1, beta_neg=0.1):
    device = next(model_pos.parameters()).device
    X = X.to(device=device, dtype=torch.float32)
    y = y.to(device=device)

    model_pos.eval()
    model_neg.eval()
    with torch.no_grad():
        s1 = model_pos(X).flatten()
        s2 = model_neg(X).flatten()

    pos = (y == 1)
    neg = (y == -1)
    n1 = int(pos.sum().item())
    n2 = int(neg.sum().item())

    n02_1 = int((((s1 >= beta_pos) & (s2 <  beta_neg)) & pos).sum().item())
    n12_1 = int((((s1 >= beta_pos) & (s2 >= beta_neg)) & pos).sum().item())

    n02_2 = int((((s2 >= beta_neg) & (s1 <  beta_pos)) & neg).sum().item())
    n12_2 = int((((s2 >= beta_neg) & (s1 >= beta_pos)) & neg).sum().item())

    a1 = n02_1 / n1 if n1 else 0.0
    a2 = n02_2 / n2 if n2 else 0.0
    F12 = 2 * a1 * a2 / (a1 + a2 + 1e-12)

    b1 = n12_1 / n1 if n1 else 0.0
    b2 = n12_2 / n2 if n2 else 0.0
    G12 = 2 * b1 * b2 / (b1 + b2 + 1e-12)

    return dict(
        n1=n1, n2=n2,
        n02_1=n02_1, n12_1=n12_1,
        n02_2=n02_2, n12_2=n12_2,
        a1=a1, a2=a2, F12=F12,
        b1=b1, b2=b2, G12=G12
    )


def pick_beta(scores_pos: torch.Tensor,
              scores_noise: torch.Tensor,
              scores_neg: Optional[torch.Tensor] = None,
              n_quantiles: int = 19) -> Tuple[float, float]:

    s_pos = scores_pos.detach().flatten()
    s_no  = scores_noise.detach().flatten()
    s_neg = scores_neg.detach().flatten() if (scores_neg is not None) else None

    if s_pos.numel() == 0:
        return 0.0, -1.0

    device = s_pos.device

    qs = torch.linspace(0.05, 0.95, steps=n_quantiles, device=device)

    if s_neg is None:
        merged = torch.cat([s_pos, s_no])
    else:
        merged = torch.cat([s_pos, s_no, s_neg])

    pool = torch.quantile(merged, qs)

    best_J, best_b = -1.0, float("nan")

    for b_tensor in pool:
        b = float(b_tensor.item())

        tpr = (s_pos >= b).float().mean().item() if s_pos.numel() else 0.0

        if s_neg is None:
            tnr = (s_no < b).float().mean().item() if s_no.numel() else 0.0
        else:
            tnr_no  = (s_no  < b).float().mean().item() if s_no.numel()  else 0.0
            tnr_neg = (s_neg < b).float().mean().item() if s_neg.numel() else 0.0
            tnr = 0.5 * (tnr_no + tnr_neg)

        J = tpr + tnr - 1.0
        if J > best_J:
            best_J, best_b = J, b

    return best_b, best_J


def choose_coord(
    X: torch.Tensor,
    y: torch.Tensor,
    d_hidden: int,
    n_hidden_layers: int,
    num_epochs: int,
    batch_size: int = 128,
    seed: int = 0,
    output: bool = False,
):

    if torch.is_tensor(y) and y.device != X.device:
        y = y.to(X.device)

    n_features = X.shape[1]


    idx_np = torch.arange(len(y)).detach().cpu().numpy()
    y_np = y.detach().cpu().numpy()

    train_idx_np, val_idx_np = train_test_split(
        idx_np, test_size=0.2, random_state=seed, stratify=y_np
    )


    train_idx = torch.as_tensor(train_idx_np, dtype=torch.long, device=X.device)
    val_idx   = torch.as_tensor(val_idx_np,   dtype=torch.long, device=X.device)


    X_tr, y_tr   = X[train_idx], y[train_idx]
    X_val, y_val = X[val_idx],   y[val_idx]

    fmin_all, fmax_all = compute_feature_bounds(TensorDataset(X_tr, y_tr))
    fmin_all = fmin_all.to(X_tr.device)
    fmax_all = fmax_all.to(X_tr.device)

    rows = []
    best_idx, best_key = None, -float("inf")

    for i in range(n_features):

        X_tr_upd  = torch.cat((X_tr[:, :i],  X_tr[:, i+1:]),  dim=1)
        X_val_upd = torch.cat((X_val[:, :i], X_val[:, i+1:]), dim=1)

        fmin_upd = torch.cat((fmin_all[:i], fmin_all[i+1:]))
        fmax_upd = torch.cat((fmax_all[:i], fmax_all[i+1:]))

        n_pos_val = int((y_val == 1).sum().item())
        n_neg_val = int((y_val == -1).sum().item())
        n_noise_for_pos = max(1, n_pos_val)
        n_noise_for_neg = max(1, n_neg_val)

        noise_pos_ds = UniformNoiseDataset(fmin_upd, fmax_upd, n_noise_for_pos, label=0, dtype=torch.float32)
        noise_neg_ds = UniformNoiseDataset(fmin_upd, fmax_upd, n_noise_for_neg, label=0, dtype=torch.float32)

        X_noise_pos = torch.stack([noise_pos_ds[j][0] for j in range(len(noise_pos_ds))], dim=0)
        X_noise_neg = torch.stack([noise_neg_ds[j][0] for j in range(len(noise_neg_ds))], dim=0)

        train_data_upd = TensorDataset(X_tr_upd, y_tr)

        model_pos, model_neg = model_train(
            train_data=train_data_upd,
            d_hidden=d_hidden,
            n_hidden_layers=n_hidden_layers,
            num_epochs=num_epochs,
            batch_size=batch_size,
            seed=seed,
            output=output,
        )

        model_pos.eval()
        model_neg.eval()
        dev_pos = next(model_pos.parameters()).device
        dev_neg = next(model_neg.parameters()).device

        with torch.no_grad():
            s1_pos = model_pos(X_val_upd[(y_val == 1)]).flatten()
            s1_no  = model_pos(X_noise_pos).flatten()
            beta_pos, _ = pick_beta(s1_pos, s1_no)

            s2_neg = model_neg(X_val_upd[(y_val == -1)]).flatten()
            s2_no  = model_neg(X_noise_neg).flatten()
            beta_neg, _ = pick_beta(s2_neg, s2_no)

        cG_grid = [0.1, 0.25, 0.5, 0.75, 1.0]
        best_S, best_tuple, best_m = -float("inf"), None, None

        for cG in cG_grid:
            m_try = metrics_from_sheet(
                X_val_upd, y_val, model_pos, model_neg,
                beta_pos=beta_pos, beta_neg=beta_neg
            )
            S_try = m_try["F12"] - cG * m_try["G12"]
            if S_try > best_S:
                best_S, best_tuple, best_m = S_try, (beta_pos, beta_neg, cG), m_try

        m = dict(best_m)
        m["removed_idx"] = i
        m["beta_pos_best"], m["beta_neg_best"], m["cG_best"] = best_tuple
        m["S"] = best_S
        rows.append(m)

        if best_S > best_key:
            best_key = best_S
            best_idx = i

    df = (
        pd.DataFrame(rows)
          .sort_values(["S", "F12", "G12"], ascending=[False, False, True])
          .reset_index(drop=True)
    )
    return best_idx, df

def set_seed(seed: int = 42):
    import random
    import numpy as np
    import torch

    random.seed(seed)

    np.random.seed(seed)

    torch.manual_seed(seed)

    if torch.cuda.is_available():
        torch.cuda.manual_seed_all(seed)


class ReducerBase:
    def fit(self, X, y=None): ...
    def transform(self, X): ...
    def fit_transform(self, X, y=None):
        self.fit(X, y)
        return self.transform(X)

    @staticmethod
    def _to_numpy(X):
        if torch.is_tensor(X):
            return X.detach().cpu().numpy()
        return np.asarray(X)

    @staticmethod
    def _to_numpy_y(y):
        if y is None:
            return None
        if torch.is_tensor(y):
            return y.detach().cpu().numpy()
        return np.asarray(y)

class MLPReducer(ReducerBase):
    def __init__(self, d_hidden, n_hidden_layers, num_epochs, batch_size, n_select, seed=0, output=False, device: str = "cpu",):
        self.d_hidden = d_hidden
        self.n_hidden_layers = n_hidden_layers
        self.num_epochs = num_epochs
        self.batch_size = batch_size
        self.seed = seed
        self.output = output
        self.n_select = n_select

        self.selected_indices_ = None
        self.removed_indices_ = None
        self.history_ = []
        self.device = torch.device(device)
        

    def _to_torch2d(self, X):
        if isinstance(X, torch.Tensor):
            assert X.ndim == 2, "X must be 2D"
            return X.to(device=self.device, dtype=torch.float32)
        X = np.asarray(X)
        assert X.ndim == 2, "X must be 2D"
        return torch.from_numpy(X).to(device=self.device, dtype=torch.float32)

    def _to_torch1d_int(self, y):
        if isinstance(y, torch.Tensor):
            assert y.ndim == 1, "y must be 1D"
            return y.to(device=self.device, dtype=torch.int64)
        y = np.asarray(y).ravel()
        return torch.from_numpy(y.astype(np.int64)).to(device=self.device)


    def fit(self, X, y):
        X_t = self._to_torch2d(X)
        y_t = self._to_torch1d_int(y)

        y_t = y_t.to(device=X_t.device)

        d = X_t.shape[1]
        keep = list(range(d))
        removed = []

        while len(keep) > self.n_select:
            X_cur = X_t[:, keep]

            best_idx, df = choose_coord(
                X=X_cur, y=y_t,
                d_hidden=self.d_hidden,
                n_hidden_layers=self.n_hidden_layers,
                num_epochs=self.num_epochs,
                batch_size=self.batch_size,
                seed=self.seed,
                output=self.output,
            )

            removed_global = keep[best_idx]
            removed.append(removed_global)
            del keep[best_idx]

            self.history_.append({
                "removed_global": removed_global,
                "n_left": len(keep),
                "df": df,
            })

        self.selected_indices_ = keep
        self.removed_indices_ = removed
        return self

    def transform(self, X):
        assert self.selected_indices_ is not None, "Call fit() before transform()."

        if isinstance(X, torch.Tensor):
            return X[:, self.selected_indices_]
        X = np.asarray(X)
        return X[:, self.selected_indices_]

    def fit_transform(self, X, y):
        return self.fit(X, y).transform(X)

    def get_support(self, indices=False):
        assert self.selected_indices_ is not None, "Call fit() first."
        if indices:
            return np.array(self.selected_indices_, dtype=int)

        if self.selected_indices_:
            mask_len = max(self.selected_indices_) + 1
        else:
            mask_len = self.n_select

        mask = np.zeros(mask_len, dtype=bool)
        if self.selected_indices_:
            mask[self.selected_indices_] = True
        return mask
\end{lstlisting}

\subsection*{ПРИЛОЖЕНИЕ Б}
\addcontentsline{toc}{subsection}{ПРИЛОЖЕНИЕ Б}

Ниже приведен код, относящийся к RaSE методу.

\begin{lstlisting}[
    language=Python,
    caption={Метода RaSE с использованием унарной классификации и SHAP},
    label={lst:process_datasets},
    captionpos=b
]

def sample_uniform(fmin: torch.Tensor, fmax: torch.Tensor, n: int) -> torch.Tensor:

    fmax = fmax.to(device=fmin.device, dtype=fmin.dtype)


    u = torch.rand(n, fmin.numel(), dtype=fmin.dtype, device=fmin.device)
    return fmin.unsqueeze(0) + (fmax - fmin).unsqueeze(0) * u

def tune_betas_joint(
    model_pos: torch.nn.Module,
    model_neg: torch.nn.Module,
    X_val: torch.Tensor,
    y_val: torch.Tensor,
    fmin: torch.Tensor,
    fmax: torch.Tensor,
    n_grid: int = 19,
    cG_grid: Iterable[float] = (0.25, 0.5, 0.75, 1.0),
) -> Tuple[float, float, float, float]:
    """
    Совместно подбирает (beta_pos, beta_neg) и cG, максимизируя
    S = F12 - cG * G12 на валидации.
    Возвращает: (beta_pos, beta_neg, best_cG, best_S).
    """

    device = next(model_pos.parameters()).device
    Xv = X_val.to(device=device, dtype=torch.float32)
    yv = y_val.to(device=device, dtype=torch.int64)

    model_pos.eval(); model_neg.eval()
    with torch.no_grad():
        s1_all = model_pos(Xv).flatten()
        s2_all = model_neg(Xv).flatten()
        pos = (yv == 1)
        neg = (yv == -1)

        n1 = max(1, int(pos.sum().item()))
        n2 = max(1, int(neg.sum().item()))
        Xno1 = sample_uniform(fmin.to(device), fmax.to(device), n1)
        Xno2 = sample_uniform(fmin.to(device), fmax.to(device), n2)
        s1_no = model_pos(Xno1).flatten()
        s2_no = model_neg(Xno2).flatten()

        q = torch.linspace(0.05, 0.95, steps=n_grid, device=device)
        pool1 = torch.quantile(torch.cat([s1_all[pos], s1_all[neg], s1_no]), q)
        pool2 = torch.quantile(torch.cat([s2_all[neg], s2_all[pos], s2_no]), q)

        n1 = max(1, int(pos.sum().item()))
        n2 = max(1, int(neg.sum().item()))

    best_S = -float("inf"); best_bp=0.0; best_bn=0.0; best_cG=1.0

    for bpos_t in pool1:
        bpos = float(bpos_t.item())
        m1 = (s1_all >= bpos)
        for bneg_t in pool2:
            bneg = float(bneg_t.item())
            m2 = (s2_all >= bneg)

            n02_1 = int(((m1 & ~m2) & pos).sum().item())
            n02_2 = int(((m2 & ~m1) & neg).sum().item())
            a1 = n02_1 / n1
            a2 = n02_2 / n2
            F12 = 2*a1*a2 / (a1+a2+1e-12)

            n12_1 = int(((m1 & m2) & pos).sum().item())
            n12_2 = int(((m2 & m1) & neg).sum().item())
            b1 = n12_1 / n1
            b2 = n12_2 / n2
            G12 = 2*b1*b2 / (b1+b2+1e-12)

            for cG in cG_grid:
                S = F12 - float(cG) * G12
                if S > best_S:
                    best_S = S; best_bp=bpos; best_bn=bneg; best_cG=float(cG)

    return best_bp, best_bn, best_cG, best_S


def _project_capped_simplex(p: np.ndarray, L: float, U: float) -> np.ndarray:
    p = np.asarray(p, dtype=float)
    n = p.size
    L = float(L); U = float(U)
    L = min(L, 1.0/n)  
    U = max(U, 1.0/n)

    x = np.clip(p, L, U)
    r = 1.0 - x.sum()
    if abs(r) < 1e-12:
        return x / x.sum()

    for _ in range(5): 
        if r > 0:
            mask = (x < U - 1e-12)
            if not np.any(mask):
                break
            w = p[mask].copy()
            w[w < 0] = 0.0
            if not np.any(w):
                w = np.ones_like(w)
            w = w / w.sum()
            delta = np.minimum(U - x[mask], r * w / (w.sum() + 1e-12))
            add = delta * (r / (delta.sum() + 1e-12))
            x[mask] += add
            r = 1.0 - x.sum()
        else:
            mask = (x > L + 1e-12)
            if not np.any(mask):
                break
            w = p[mask].copy()
            w[w < 0] = 0.0
            if not np.any(w):
                w = np.ones_like(w)
            w = w / w.sum()
            delta = np.minimum(x[mask] - L, (-r) * w / (w.sum() + 1e-12))
            sub = delta * ((-r) / (delta.sum() + 1e-12))
            x[mask] -= sub
            r = 1.0 - x.sum()

        if abs(r) < 1e-12:
            break

    x = np.clip(x, L, U)
    x /= x.sum()
    return x

def dynamic_caps(n: int, q: int,
                 alpha_top: float = 0.80,
                 kappa: float = 0.50,
                 u_clip: float = 0.25,
                 eps: float = 1e-6) -> tuple[float, float]:

    assert n >= 1 and 1 <= q <= n, "dynamic_caps: bad n or q"
    U = min(alpha_top / max(q, 1), u_clip)
    U = max(U, 1.0 / n + eps)  

    if n > q:
        L_max = max(0.0, (1.0 - q * U) / (n - q))
    else:
        L_max = 0.0

    L = kappa * L_max
    L = max(0.0, min(L, U - eps))
    return L, U


class EnsembleReducer(ReducerBase):
    def __init__(
        self,
        d_hidden,
        n_hidden_layers,
        num_epochs,
        batch_size,
        n_select,
        num_models,
        num_attempts,
        num_coords,
        num_samples,
        seed=0,
        output=False,
        device: str = "cpu",
        use_gpu_shap: bool = False
    ):
        self.d_hidden = d_hidden
        self.n_hidden_layers = n_hidden_layers
        self.num_epochs = num_epochs
        self.batch_size = batch_size
        self.seed = seed
        self.output = output
        self.n_select = n_select
        self.num_samples = num_samples
        self.num_models = num_models
        self.num_attempts = num_attempts
        self.num_coords = num_coords

        self.list_model_pos = []
        self.list_model_neg = []
        self.list_model_subs = []
        self.list_beta_pos = []
        self.list_beta_neg = []
        self.list_best_cG = []
        self.list_best_S = []
        self.coord_prob = None
        self.history_ = []
        self.rng = np.random.default_rng(self.seed)

        self.selected_indices_ = None

        self.device = torch.device(device)
        self.use_gpu_shap = use_gpu_shap

    def _to_torch2d(self, X):
        if isinstance(X, torch.Tensor):
            assert X.ndim == 2, "X must be 2D"
            return X.to(dtype=torch.float32)
        X = np.asarray(X)
        assert X.ndim == 2, "X must be 2D"
        return torch.from_numpy(X).to(dtype=torch.float32)

    def _to_torch1d_int(self, y):
        if isinstance(y, torch.Tensor):
            assert y.ndim == 1, "y must be 1D"
            return y.to(dtype=torch.int64)
        y = np.asarray(y).ravel()
        return torch.from_numpy(y.astype(np.int64))

    @staticmethod
    def _shap_mean_pos(expl, X):
        vals = expl.shap_values(X)
        if isinstance(vals, list):
            vals = np.array(vals)
            if vals.ndim == 3: 
                vals = vals.mean(axis=0)
        V = np.atleast_2d(np.asarray(vals, float))
        return np.maximum(V, 0.0).mean(axis=0)

    @staticmethod
    def _robust_norm(v: np.ndarray) -> np.ndarray:
        v = np.asarray(v, dtype=float)
        v = np.maximum(v, 0.0)
        if not np.any(np.isfinite(v)):
            return np.full_like(v, 1.0 / max(1, v.size))
        p95 = np.percentile(v[np.isfinite(v)], 95.0)
        if p95 > 0:
            v = np.minimum(v, p95)
        s = v.sum()
        return (v / s) if (np.isfinite(s) and s > 0) else np.full_like(v, 1.0 / max(1, v.size))


    def fit(self, X, y):

        X_t = self._to_torch2d(X).to(self.device)
        y_t = self._to_torch1d_int(y).to(self.device)

        n_feat = X_t.shape[1]
        n_samples_total = X_t.shape[0]

    
        y_np = y_t.detach().cpu().numpy().ravel()
        idx_np = np.arange(n_samples_total)

        train_idx_np, val_idx_np = train_test_split(
            idx_np,
            test_size=0.2,
            random_state=self.seed,
            stratify=y_np
        )

        train_idx = torch.as_tensor(train_idx_np, dtype=torch.long, device=X_t.device)
        val_idx   = torch.as_tensor(val_idx_np,   dtype=torch.long, device=X_t.device)

        X_tr, y_tr   = X_t[train_idx], y_t[train_idx]
        X_val, y_val = X_t[val_idx],   y_t[val_idx]


        fmin_all, fmax_all = compute_feature_bounds(TensorDataset(X_tr, y_tr))
        fmin_all = fmin_all.to(device=X_tr.device, dtype=X_tr.dtype)
        fmax_all = fmax_all.to(device=X_tr.device, dtype=X_tr.dtype)

        self.coord_prob = np.full(n_feat, 1.0 / n_feat, dtype=float)

        for k in range(self.num_samples):

            for i in range(self.num_models):
                best_model_pos = best_model_neg = None
                best_subs_idx = None
                best_beta_pos = best_beta_neg = None
                best_best_S = -float('inf')
                best_best_cG = None

                for j in range(self.num_attempts):
                    lam = 0.3
                    if self.rng.random() < lam:
                        subs_idx = self.rng.choice(n_feat, size=self.num_coords, replace=False)
                    else:
                        subs_idx = self.rng.choice(
                            n_feat, size=self.num_coords, replace=False, p=self.coord_prob
                        )

                    subs_idx = np.sort(subs_idx)
                    subs_idx_t = torch.as_tensor(subs_idx, dtype=torch.long, device=X_tr.device)

                    X_sub_tr  = X_tr[:, subs_idx_t]
                    X_sub_val = X_val[:, subs_idx_t]

                    fmin_upd = fmin_all[subs_idx_t]
                    fmax_upd = fmax_all[subs_idx_t]

                    train_ds = TensorDataset(X_sub_tr, y_tr)

                    model_pos, model_neg = model_train(
                        d_hidden=self.d_hidden,
                        n_hidden_layers=self.n_hidden_layers,
                        train_data=train_ds,
                        num_epochs=self.num_epochs,
                        batch_size=self.batch_size,
                        seed=self.seed + 1337 * k + j,
                        output=self.output,
                        device=self.device,
                        feature_min=fmin_upd,
                        feature_max=fmax_upd,
                    )
                    model_pos.eval()
                    model_neg.eval()

                    bpos, bneg, cG_star, S_star = tune_betas_joint(
                        model_pos, model_neg,
                        X_sub_val, y_val,
                        fmin_upd, fmax_upd,
                        n_grid=19,
                        cG_grid=(0.25, 0.5, 0.75, 1.0)
                    )

                    if S_star > best_best_S:
                        best_best_S   = S_star
                        best_subs_idx = subs_idx
                        best_model_pos = model_pos
                        best_model_neg = model_neg
                        best_beta_pos  = bpos
                        best_beta_neg  = bneg
                        best_best_cG   = cG_star

                self.list_model_pos.append(best_model_pos)
                self.list_model_neg.append(best_model_neg)
                self.list_model_subs.append(best_subs_idx)
                self.list_beta_pos.append(best_beta_pos)
                self.list_beta_neg.append(best_beta_neg)
                self.list_best_cG.append(best_best_cG)
                self.list_best_S.append(best_best_S)

            start = len(self.list_model_pos) - self.num_models
            end   = len(self.list_model_pos)
            S_slice = np.asarray(self.list_best_S[start:end], dtype=float)

            keep_top = 2
            keep_top = min(keep_top, max(1, end - start))
            top_local_idx = np.argpartition(-S_slice, kth=min(keep_top, S_slice.size) - 1)[:keep_top]
            top_global_idx = (top_local_idx + start).tolist()


            coord_new_prob = np.zeros(n_feat, dtype=float)

            bg_size = min(256, X_tr.shape[0])
            ev_size = min(256, X_val.shape[0])

            ytr_np = y_tr.detach().cpu().numpy().ravel()
            yvl_np = y_val.detach().cpu().numpy().ravel()
            pos_tr = np.flatnonzero(ytr_np ==  1)
            neg_tr = np.flatnonzero(ytr_np == -1)
            pos_vl = np.flatnonzero(yvl_np ==  1)
            neg_vl = np.flatnonzero(yvl_np == -1)

            def _choose_5050(pos_idx, neg_idx, total, seed):
                rng = np.random.default_rng(seed)
                if (len(pos_idx) == 0) and (len(neg_idx) == 0):
                    return np.array([], dtype=int)
                if (len(pos_idx) == 0) or (len(neg_idx) == 0):
                    pool = np.concatenate([pos_idx, neg_idx])
                    total_eff = min(total, len(pool))
                    return rng.choice(pool, size=total_eff, replace=False)

                n_pos = total // 2
                n_neg = total - n_pos
                n_pos = min(n_pos, len(pos_idx))
                n_neg = min(n_neg, len(neg_idx))
                need  = total - (n_pos + n_neg)

                if need > 0:
                    rem_pos = len(pos_idx) - n_pos
                    add = min(need, max(0, rem_pos))
                    n_pos += add
                    need  -= add
                if need > 0:
                    rem_neg = len(neg_idx) - n_neg
                    add = min(need, max(0, rem_neg))
                    n_neg += add
                    need  -= add

                chosen_pos = (np.random.default_rng(seed).choice(pos_idx, size=n_pos, replace=False)
                              if n_pos > 0 else np.array([], dtype=int))
                chosen_neg = (np.random.default_rng(seed + 1).choice(neg_idx, size=n_neg, replace=False)
                              if n_neg > 0 else np.array([], dtype=int))
                chosen = np.concatenate([chosen_pos, chosen_neg])
                rng.shuffle(chosen)
                return chosen

            bg_idx_np = _choose_5050(pos_tr, neg_tr, bg_size, seed=self.seed + k)
            ev_idx_np = _choose_5050(pos_vl, neg_vl, ev_size, seed=self.seed + k + 1)

            if bg_idx_np.size == 0:
                bg_idx_np = np.arange(X_tr.shape[0])
            if ev_idx_np.size == 0:
                ev_idx_np = np.arange(X_val.shape[0])

            bg_idx = torch.as_tensor(bg_idx_np, dtype=torch.long, device=X_tr.device)
            ev_idx = torch.as_tensor(ev_idx_np, dtype=torch.long, device=X_val.device)

            for i_top in top_global_idx:
                subs = self.list_model_subs[i_top]
                subs_t = torch.as_tensor(subs, dtype=torch.long, device=X_tr.device)

                X_bg   = X_tr[bg_idx][:, subs_t]
                X_eval = X_val[ev_idx][:, subs_t]
                y_eval = y_val[ev_idx]

                X_pos = X_eval[(y_eval ==  1)]
                X_neg = X_eval[(y_eval == -1)]
                if X_pos.shape[0] == 0:
                    X_pos = X_eval
                if X_neg.shape[0] == 0:
                    X_neg = X_eval

                if self.use_gpu_shap and self.device.type == "cuda":
                    mp = copy.deepcopy(self.list_model_pos[i_top]).to(self.device).eval()
                    mn = copy.deepcopy(self.list_model_neg[i_top]).to(self.device).eval()

                    X_bg_expl  = X_bg  
                    X_pos_expl = X_pos
                    X_neg_expl = X_neg

                else:

                    mp = copy.deepcopy(self.list_model_pos[i_top]).to("cpu").eval()
                    mn = copy.deepcopy(self.list_model_neg[i_top]).to("cpu").eval()

                    X_bg_expl  = X_bg.detach().cpu()
                    X_pos_expl = X_pos.detach().cpu()
                    X_neg_expl = X_neg.detach().cpu()


                expl_pos = shap.GradientExplainer(mp, X_bg_expl)
                expl_neg = shap.GradientExplainer(mn, X_bg_expl)

                mu_pos = self._shap_mean_pos(expl_pos, X_pos_expl)
                mu_neg = self._shap_mean_pos(expl_neg, X_neg_expl)

                mu_pos = self._robust_norm(mu_pos)
                mu_neg = self._robust_norm(mu_neg)

                mu_pos = np.squeeze(np.asarray(mu_pos, dtype=float)).reshape(-1)
                mu_neg = np.squeeze(np.asarray(mu_neg, dtype=float)).reshape(-1)
                assert mu_pos.shape == mu_neg.shape, f"Shape mismatch: {mu_pos.shape} vs {mu_neg.shape}"

                imp_sub = 2.0 * mu_pos * mu_neg / (mu_pos + mu_neg + 1e-12)

                s = imp_sub.sum()
                if s > 0 and np.isfinite(s):
                    imp_sub /= s
                coord_new_prob[subs] += imp_sub

            tot = coord_new_prob.sum()
            if not np.isfinite(tot) or tot <= 0:
                coord_new_prob[:] = 1.0 / n_feat
            else:
                coord_new_prob /= tot

            alpha = 0.35
            coord_smooth = (1.0 - alpha) * self.coord_prob + alpha * coord_new_prob

            L, U = dynamic_caps(n_feat, self.n_select, alpha_top=0.80, kappa=0.50, u_clip=0.25)
            coord_smooth = _project_capped_simplex(coord_smooth, L=L, U=U)

            self.coord_prob = coord_smooth

            if self.output:
                print("pi:", np.round(self.coord_prob, 4))


        return self

    def coord_prob_get(self):
        return self.coord_prob

    def transform(self, X):
        assert self.coord_prob is not None, "Call fit() before transform()."
        assert hasattr(self, "n_select"), "n_select must be set in __init__()."

        k = min(self.n_select, len(self.coord_prob))
        topk_idx = np.argsort(self.coord_prob)[-k:]
        topk_idx = np.sort(topk_idx)
        self.selected_indices_ = topk_idx

        if isinstance(X, torch.Tensor):
            assert X.ndim == 2, "X must be 2D tensor"
            return X[:, topk_idx]
        else:
            X = np.asarray(X)
            assert X.ndim == 2, "X must be 2D array"
            return X[:, topk_idx]


\end{lstlisting}

\subsection*{ПРИЛОЖЕНИЕ В}
\addcontentsline{toc}{subsection}{ПРИЛОЖЕНИЕ В}

Ниже приведен код относящийся к проведению экспериментов по сравнению методов с существующими.

\begin{lstlisting}[
    language=Python,
    caption={Проводимые эксперименты},
    label={lst:process_datasets},
    captionpos=b
]

def make_dataset_A(seed=42):
    X, y = make_classification(
        n_samples=3000, n_features=8, n_informative=3, n_redundant=2,
        class_sep=1.2, flip_y=0.02, random_state=seed)
    X[:, 0] = np.sin(X[:, 0]) + 0.1 * X[:, 1]
    X[:, 2] = X[:, 2] ** 2
    y = y.astype(int)
    y = 2*y - 1
    return X.astype(np.float64), y.astype(int)

def make_dataset_B(seed=42):
    X, y = make_classification(
        n_samples=6000, n_features=30, n_informative=5, n_redundant=5,
        class_sep=0.8, flip_y=0.01, weights=[0.75, 0.25], random_state=seed)
    X[:, 10] = np.sin(X[:, 0]) * X[:, 1]
    X[:, 11] = X[:, 2] ** 2
    X[:, 12] = np.exp(-X[:, 3]**2)
    X[:, 20:25] = X[:, :5] + 0.05 * np.random.randn(len(X), 5)
    y = y.astype(int)
    y = 2*y - 1
    return X.astype(np.float64), y.astype(int)

def make_dataset_A_real(seed: int = 42) -> tuple[np.ndarray, np.ndarray]:
    ds = fetch_ucirepo(id=545) 
    X_df = ds.data.features
    y_df = ds.data.targets

    X = X_df.to_numpy(dtype=np.float64)
    y_raw = np.asarray(y_df.squeeze()) 

    classes = np.unique(y_raw)

    pos_label = "Osmancik" if "Osmancik" in classes else classes[-1]
    y = np.where(y_raw == pos_label, 1, -1).astype(int)

    return X.astype(np.float64), y.astype(int)

def make_dataset_B_real(seed: int = 42) -> tuple[np.ndarray, np.ndarray]:
    ds = fetch_ucirepo(id=327) 
    X = ds.data.features.to_numpy(dtype=np.float64)
    y_raw = np.asarray(ds.data.targets.squeeze())

    y_num = y_raw.astype(int)

    y = np.where(y_num > 0, 1, -1).astype(int)

    return X.astype(np.float64), y.astype(int)

class PCAReducer(ReducerBase):
    def __init__(self, n_components):
        self.pca = PCA(n_components=n_components, random_state=42)

    def fit(self, X, y=None):
        Xn = self._to_numpy(X)
        self.pca.fit(Xn)
        return self

    def transform(self, X):
        Xn = self._to_numpy(X)
        return self.pca.transform(Xn)

class PLSReducer(ReducerBase):
    def __init__(self, n_components):
        self.pls = PLSRegression(n_components=n_components, scale=False)

    def fit(self, X, y):
        Xn = self._to_numpy(X)
        yn = self._to_numpy_y(y).ravel().astype(float)
        self.pls.fit(Xn, yn)
        return self

    def transform(self, X):
        Xn = self._to_numpy(X)
        return self.pls.transform(Xn)


class UMAPReducer(ReducerBase):
    def __init__(self, n_components):
        self.umap = UMAP(
            n_components=n_components,
            random_state=42,
            n_neighbors=15,
            min_dist=0.1,
            target_metric='categorical'
        )

    def fit(self, X, y):
        Xn = self._to_numpy(X)
        yn = self._to_numpy_y(y)
        self.umap.fit(Xn, yn)
        return self

    def transform(self, X):
        Xn = self._to_numpy(X)
        return self.umap.transform(Xn)

class HSICSelector(ReducerBase):
    def __init__(self, n_select):
        self.n_select = int(n_select)
        self.model = HSICLasso()
        self.mask_ = None

    def fit(self, X, y):
        Xn = self._to_numpy(X)
        yn = self._to_numpy_y(y).ravel().astype(int)

        self.model.input(Xn, yn)
        self.model.classification(self.n_select)

        idx = np.array(self.model.get_index(), dtype=int)
        self.mask_ = np.zeros(Xn.shape[1], dtype=bool)
        self.mask_[idx] = True
        return self

    def transform(self, X):
        Xn = self._to_numpy(X)
        return Xn[:, self.mask_]

def get_models():
    return {
        "MLP": MLPClassifier(hidden_layer_sizes=(64,32), activation='relu',
                             alpha=1e-4, batch_size=256, learning_rate_init=1e-3,
                             max_iter=200, early_stopping=True, n_iter_no_change=10, random_state=42),
        "HGBT": HistGradientBoostingClassifier(random_state=42)
    }


def evaluate_reduction(
    X: np.ndarray,
    y: np.ndarray,
    reducer_name: str,
    reducer_ctor: Any,
    q: int,
    dataset_name: str,
    n_splits: int = 5,
    reducer_kwargs: Dict[str, Any] = None,
    compute_unary_on_reduced: bool = True,
    unary_params: Dict[str, Any] = None,
) -> Tuple[List[Dict[str, Any]], List[Dict[str, Any]]]:


    reducer_kwargs = reducer_kwargs or {}
    unary_params = unary_params or dict(
        d_hidden=48,
        n_hidden_layers=2,
        num_epochs=60,
        batch_size=256,
        beta_grid=31,
        cG_grid=(0.10, 0.20, 0.25, 0.33, 0.50, 0.67, 0.75, 1.00, 1.25),

    )

    y_orig = np.asarray(y).ravel()
    if set(np.unique(y_orig)) == {-1, 1}:
        y_bin = (y_orig == 1).astype(int)
    else:
        y_bin = y_orig.astype(int)

    skf = StratifiedKFold(n_splits=n_splits, shuffle=True, random_state=42)
    rows_cls: List[Dict[str, Any]] = []
    rows_unary: List[Dict[str, Any]] = []

    for fold, (tr, te) in enumerate(skf.split(X, y_bin), 1):
        Xtr, Xte = X[tr], X[te]
        ytr_bin, yte_bin = y_bin[tr], y_bin[te]
        ytr_orig, yte_orig = y_orig[tr], y_orig[te]


        scaler = StandardScaler()
        Xtr_s = scaler.fit_transform(Xtr)
        Xte_s = scaler.transform(Xte)


        base_scores = {}
        models = get_models()
        for mname, base_model in models.items():
            m0 = base_model.__class__(**base_model.get_params())
            t0 = perf_counter()
            m0.fit(Xtr_s, ytr_bin)
            t0 = perf_counter() - t0
            p0 = m0.predict_proba(Xte_s)[:, 1] if hasattr(m0, "predict_proba") else m0.decision_function(Xte_s)
            auc0 = roc_auc_score(yte_bin, p0)
            ap0 = average_precision_score(yte_bin, p0)
            base_scores[mname] = (auc0, ap0, t0)


        red = reducer_ctor(q=q, **reducer_kwargs)
        tR = perf_counter()
        Ztr = red.fit_transform(Xtr_s, ytr_orig)
        Zte = red.transform(Xte_s)
        tR = perf_counter() - tR


        for mname, base_model in models.items():
            auc0, ap0, t0 = base_scores[mname]
            m1 = base_model.__class__(**base_model.get_params())
            t1 = perf_counter()
            m1.fit(Ztr, ytr_bin)
            t1 = perf_counter() - t1
            p1 = m1.predict_proba(Zte)[:, 1] if hasattr(m1, "predict_proba") else m1.decision_function(Zte)
            auc1 = roc_auc_score(yte_bin, p1)
            ap1 = average_precision_score(yte_bin, p1)

            rows_cls.append({
                "dataset": dataset_name, "method": reducer_name, "q": q, "model": mname, "fold": fold,
                "AUC": auc1, "PR_AUC": ap1, "delta_AUC": auc1 - auc0, "delta_PR_AUC": ap1 - ap0,
                "t_reducer": tR, "t_model": t1, "t_model_base": t0
            })


        if compute_unary_on_reduced:
            dev_param = unary_params.get("device", None)
            if dev_param is not None:
                unary_device_preferred = torch.device(dev_param)
            else:
                if torch.cuda.is_available():
                    unary_device_preferred = torch.device("cuda")
                elif torch.backends.mps.is_available():
                    unary_device_preferred = torch.device("mps")
                else:
                    unary_device_preferred = torch.device("cpu")

            ytr_pm = np.where(ytr_orig == 1, 1, -1)
            yte_pm = np.where(yte_orig == 1, 1, -1)

            sss = StratifiedShuffleSplit(n_splits=1, test_size=0.2, random_state=42)
            (fit_idx, val_idx), = sss.split(Ztr, (ytr_pm == 1).astype(int))
            Ztr_fit, Ztr_val = Ztr[fit_idx], Ztr[val_idx]
            y_fit_pm, y_val_pm = ytr_pm[fit_idx], ytr_pm[val_idx]


            X_fit_t = torch.as_tensor(Ztr_fit, dtype=torch.float32, device=unary_device_preferred)
            y_fit_t = torch.as_tensor(y_fit_pm, dtype=torch.int64, device=unary_device_preferred)
            train_ds_fit = TensorDataset(X_fit_t, y_fit_t)

            model_pos, model_neg = model_train(
                d_hidden=unary_params.get("d_hidden", 32),
                n_hidden_layers=unary_params.get("n_hidden_layers", 1),
                train_data=train_ds_fit,
                num_epochs=unary_params.get("num_epochs", 50),
                batch_size=unary_params.get("batch_size", 256),
                seed=unary_params.get("seed", 0),
                output=False,
            )
            model_pos.eval()
            model_neg.eval()

            unary_device = next(model_pos.parameters()).device 

            X_val_t = torch.as_tensor(Ztr_val, dtype=torch.float32, device=unary_device)
            y_val_t = torch.as_tensor(y_val_pm, dtype=torch.int64, device=unary_device)

            fmin, fmax = compute_feature_bounds(TensorDataset(X_fit_t, y_fit_t))
            fmin = fmin.to(unary_device)
            fmax = fmax.to(unary_device)

            beta_pos, beta_neg, cG_star, S_val = tune_betas_joint(
                model_pos, model_neg,
                X_val_t, y_val_t,
                fmin, fmax,
                n_grid=unary_params.get("beta_grid", 19),
                cG_grid=unary_params.get("cG_grid", (0.25, 0.5, 0.75, 1.0))
            )


            X_te_t = torch.as_tensor(Zte, dtype=torch.float32, device=unary_device)
            yte_pm_t = torch.as_tensor(yte_pm, dtype=torch.int64, device=unary_device)

            m_un = metrics_from_sheet(
                X_te_t, yte_pm_t,
                model_pos, model_neg,
                beta_pos=float(beta_pos), beta_neg=float(beta_neg)
            )


            with torch.no_grad():
                s1_te = model_pos(X_te_t).flatten()
                s2_te = model_neg(X_te_t).flatten()

            pos_hit = (s1_te >= float(beta_pos)) & (s2_te <  float(beta_neg))
            neg_hit = (s2_te >= float(beta_neg)) & (s1_te <  float(beta_pos))
            coverage = float((pos_hit | neg_hit).float().mean().item())

            S_test = float(m_un["F12"] - cG_star * m_un["G12"])

            rows_unary.append({
                "dataset": dataset_name, "method": reducer_name, "q": q, "model": "UnaryMLP", "fold": fold,
                "F12": float(m_un["F12"]), "G12": float(m_un["G12"]),
                "a1": float(m_un["a1"]), "a2": float(m_un["a2"]),
                "b1": float(m_un["b1"]), "b2": float(m_un["b2"]),
                "n1": int(m_un["n1"]), "n2": int(m_un["n2"]),
                "coverage": coverage,
                "beta_pos": float(beta_pos), "beta_neg": float(beta_neg),
                "cG": float(cG_star), "S_val": float(S_val), "S_test": S_test,
                "t_reducer": tR, "t_model": 0.0,
                "unary_device": str(unary_device),
            })

    return rows_cls, rows_unary

def run_all_A(
    include_unary: bool = False,
    unary_params: dict | None = None,
) -> Tuple[pd.DataFrame, pd.DataFrame]:
    configs = [
        ("DS-A", make_dataset_A, [2, 3, 5]),
        # ("DS-B", make_dataset_B, [5, 10, 15]),
    ]

    methods = [
        ("PCA",        lambda q, **_: PCAReducer(q)),
        ("PLS",        lambda q, **_: PLSReducer(q)),
        ("UMAP_sup",   lambda q, **_: UMAPReducer(q)),
        ("HSIC_Lasso", lambda q, **_: HSICSelector(q)),
        ("MLP_unar",   lambda q, **kw: MLPReducer(
            d_hidden=64,
            n_hidden_layers=2,
            num_epochs=40,
            batch_size=256,
            n_select=q,
            seed=0,
            output=False,
            **kw,   
        )),
    ]


    if torch.cuda.is_available():
        DEVICE = torch.device("cuda")
    elif torch.backends.mps.is_available():
        DEVICE = torch.device("mps")
    else:
        DEVICE = torch.device("cpu")

    print("Using device:", DEVICE)
    if DEVICE.type == "cuda":
        print("CUDA name:", torch.cuda.get_device_name(0))

    base_unary_params = dict(
        d_hidden=32,
        n_hidden_layers=1,
        num_epochs=50,
        batch_size=256,
    )
    if unary_params is not None:
        base_unary_params.update(unary_params)
    base_unary_params["device"] = str(DEVICE)  

    rows_cls_all: List[Dict[str, Any]] = []

    for dname, maker, qgrid in configs:
        X, y = maker()
        for name, ctor in methods:
            for q in qgrid:
                if name == "MLP_unar":
                    reducer_kwargs = {
                        "device": str(DEVICE),  
                    }
                else:
                    reducer_kwargs = {}

                cls_rows, _un_rows = evaluate_reduction(
                    X, y,
                    reducer_name=name,
                    reducer_ctor=ctor,
                    q=q,
                    dataset_name=dname,
                    reducer_kwargs=reducer_kwargs,
                    compute_unary_on_reduced=include_unary,
                    unary_params=base_unary_params,
                )
                rows_cls_all.extend(cls_rows)


    df = pd.DataFrame(rows_cls_all)

    for col in ["AUC", "PR_AUC", "delta_AUC", "delta_PR_AUC",
                "t_reducer", "t_model", "t_model_base"]:
        if col in df.columns:
            df[col] = pd.to_numeric(df[col], errors="coerce")

    agg = (
        df.groupby(["dataset", "method", "q", "model"], dropna=False)
          .agg(
              AUC_mean=("AUC", "mean"),
              PR_AUC_mean=("PR_AUC", "mean"),
              dAUC_mean=("delta_AUC", "mean"),
              dPR_mean=("delta_PR_AUC", "mean"),
              t_red_med=("t_reducer", "median"),
              t_mod_med=("t_model", "median"),
          )
          .reset_index()
          .sort_values(["dataset", "model", "dAUC_mean"], ascending=[True, True, False])
    )

    return df, agg



def run_all_B(
    include_unary: bool = True,
    unary_params: dict | None = None,
) -> Tuple[pd.DataFrame, pd.DataFrame]:
    configs = [
        ("DS-B", make_dataset_B, [5, 10, 15]),
    ]

    methods = [
        ("PCA",        lambda q, **_: PCAReducer(q)),
        ("PLS",        lambda q, **_: PLSReducer(q)),
        ("UMAP_sup",   lambda q, **_: UMAPReducer(q)),
        ("HSIC_Lasso", lambda q, **_: HSICSelector(q)),
        ("Ensembleunar", lambda q, **kw: EnsembleReducer(
            d_hidden=32, n_hidden_layers=2,
            num_epochs=8, batch_size=1024,
            n_select=q,
            num_models=4,
            num_attempts=20,
            num_coords=4,
            num_samples=10,
            seed=42, output=True,
            **kw,
        )),
    ]

    if torch.cuda.is_available():
        DEVICE = torch.device("cuda")
    elif torch.backends.mps.is_available():
        DEVICE = torch.device("mps")
    else:
        DEVICE = torch.device("cpu")

    print("Using device:", DEVICE)
    if DEVICE.type == "cuda":
        print("CUDA name:", torch.cuda.get_device_name(0))

    base_unary_params = dict(
        d_hidden=32,
        n_hidden_layers=1,
        num_epochs=50,
        batch_size=256,
    )
    if unary_params is not None:
        base_unary_params.update(unary_params)
    base_unary_params["device"] = str(DEVICE)

    rows_cls_all: List[Dict[str, Any]] = []

    for dname, maker, qgrid in configs:
        X, y = maker()
        for name, ctor in methods:
            for q in qgrid:

                if name == "Ensembleunar":
                    reducer_kwargs = {
                        "device": str(DEVICE), 
                        "use_gpu_shap": True,
                    }

                else:
                    reducer_kwargs = {}

                cls_rows, _un_rows = evaluate_reduction(
                    X, y,
                    reducer_name=name,
                    reducer_ctor=ctor,
                    q=q,
                    dataset_name=dname,
                    reducer_kwargs=reducer_kwargs, 
                    compute_unary_on_reduced=include_unary,
                    unary_params=base_unary_params,
                )
                rows_cls_all.extend(cls_rows)


    df = pd.DataFrame(rows_cls_all)

    for col in ["AUC", "PR_AUC", "delta_AUC", "delta_PR_AUC",
                "t_reducer", "t_model", "t_model_base"]:
        if col in df.columns:
            df[col] = pd.to_numeric(df[col], errors="coerce")

    agg = (
        df.groupby(["dataset", "method", "q", "model"], dropna=False)
          .agg(
              AUC_mean=("AUC", "mean"),
              PR_AUC_mean=("PR_AUC", "mean"),
              dAUC_mean=("delta_AUC", "mean"),
              dPR_mean=("delta_PR_AUC", "mean"),
              t_red_med=("t_reducer", "median"),
              t_mod_med=("t_model", "median"),
          )
          .reset_index()
          .sort_values(["dataset", "model", "dAUC_mean"], ascending=[True, True, False])
    )

    results = {
        "df": df,
        "agg": agg,
    }
    with open("results_B", "wb") as f:
        pickle.dump(results, f)

    return df, agg

def run_all_A_real(
    include_unary: bool = False,
    unary_params: dict | None = None,
) -> Tuple[pd.DataFrame, pd.DataFrame]:

    configs = [
        ("DS-A-real", make_dataset_A_real, [2, 3, 5]),
    ]

    methods = [
        ("PCA",        lambda q, **_: PCAReducer(q)),
        ("PLS",        lambda q, **_: PLSReducer(q)),
        ("UMAP_sup",   lambda q, **_: UMAPReducer(q)),
        ("HSIC_Lasso", lambda q, **_: HSICSelector(q)),
        ("MLP_unar",   lambda q, **kw: MLPReducer(
            d_hidden=64,
            n_hidden_layers=2,
            num_epochs=40,
            batch_size=256,
            n_select=q,
            seed=0,
            output=False,
            **kw,
        )),
    ]

    if torch.cuda.is_available():
        DEVICE = torch.device("cuda")
    elif torch.backends.mps.is_available():
        DEVICE = torch.device("mps")
    else:
        DEVICE = torch.device("cpu")

    print("Using device:", DEVICE)
    if DEVICE.type == "cuda":
        print("CUDA name:", torch.cuda.get_device_name(0))

    base_unary_params = dict(
        d_hidden=32,
        n_hidden_layers=1,
        num_epochs=50,
        batch_size=256,
    )
    if unary_params is not None:
        base_unary_params.update(unary_params)
    base_unary_params["device"] = str(DEVICE) 

    rows_cls_all: List[Dict[str, Any]] = []

    for dname, maker, qgrid in configs:
        X, y = maker()  
        for name, ctor in methods:
            for q in qgrid:
                if name == "MLP_unar":
                    reducer_kwargs = {"device": str(DEVICE)} 
                else:
                    reducer_kwargs = {}

                cls_rows, _un_rows = evaluate_reduction(
                    X, y,
                    reducer_name=name,
                    reducer_ctor=ctor,
                    q=q,
                    dataset_name=dname,
                    reducer_kwargs=reducer_kwargs,
                    compute_unary_on_reduced=include_unary,
                    unary_params=base_unary_params,
                )
                rows_cls_all.extend(cls_rows)

    df = pd.DataFrame(rows_cls_all)

    for col in ["AUC", "PR_AUC", "delta_AUC", "delta_PR_AUC",
                "t_reducer", "t_model", "t_model_base"]:
        if col in df.columns:
            df[col] = pd.to_numeric(df[col], errors="coerce")

    agg = (
        df.groupby(["dataset", "method", "q", "model"], dropna=False)
          .agg(
              AUC_mean=("AUC", "mean"),
              PR_AUC_mean=("PR_AUC", "mean"),
              dAUC_mean=("delta_AUC", "mean"),
              dPR_mean=("delta_PR_AUC", "mean"),
              t_red_med=("t_reducer", "median"),
              t_mod_med=("t_model", "median"),
          )
          .reset_index()
          .sort_values(["dataset", "model", "dAUC_mean"], ascending=[True, True, False])
    )

    return df, agg

def run_all_B_real(
    include_unary: bool = True,
    unary_params: dict | None = None,
) -> Tuple[pd.DataFrame, pd.DataFrame]:

    configs = [
        ("DS-B-real", make_dataset_B_real, [5, 10, 15]),
    ]

    methods = [
        ("PCA",        lambda q, **_: PCAReducer(q)),
        ("PLS",        lambda q, **_: PLSReducer(q)),
        ("UMAP_sup",   lambda q, **_: UMAPReducer(q)),
        ("HSIC_Lasso", lambda q, **_: HSICSelector(q)),
        ("Ensembleunar", lambda q, **kw: EnsembleReducer(
            d_hidden=32,
            n_hidden_layers=2,
            num_epochs=8,
            batch_size=1024,
            n_select=q,
            num_models=4,
            num_attempts=20,
            num_coords=4,
            num_samples=10,
            seed=42,
            output=True,
            **kw,
        )),
    ]

    if torch.cuda.is_available():
        DEVICE = torch.device("cuda")
    elif torch.backends.mps.is_available():
        DEVICE = torch.device("mps")
    else:
        DEVICE = torch.device("cpu")

    print("Using device:", DEVICE)
    if DEVICE.type == "cuda":
        print("CUDA name:", torch.cuda.get_device_name(0))

    base_unary_params = dict(
        d_hidden=32,
        n_hidden_layers=1,
        num_epochs=50,
        batch_size=256,
    )
    if unary_params is not None:
        base_unary_params.update(unary_params)
    base_unary_params["device"] = str(DEVICE)

    rows_cls_all: List[Dict[str, Any]] = []

    for dname, maker, qgrid in configs:
        X, y = maker() 
        for name, ctor in methods:
            for q in qgrid:
                if name == "Ensembleunar":
                    reducer_kwargs = {
                        "device": str(DEVICE), 
                        "use_gpu_shap": True,
                    }
                else:
                    reducer_kwargs = {}

                cls_rows, _un_rows = evaluate_reduction(
                    X, y,
                    reducer_name=name,
                    reducer_ctor=ctor,
                    q=q,
                    dataset_name=dname,
                    reducer_kwargs=reducer_kwargs,
                    compute_unary_on_reduced=include_unary,
                    unary_params=base_unary_params,
                )
                rows_cls_all.extend(cls_rows)

    df = pd.DataFrame(rows_cls_all)

    for col in [
        "AUC", "PR_AUC", "delta_AUC", "delta_PR_AUC",
        "t_reducer", "t_model", "t_model_base"
    ]:
        if col in df.columns:
            df[col] = pd.to_numeric(df[col], errors="coerce")

    agg = (
        df.groupby(["dataset", "method", "q", "model"], dropna=False)
          .agg(
              AUC_mean=("AUC", "mean"),
              PR_AUC_mean=("PR_AUC", "mean"),
              dAUC_mean=("delta_AUC", "mean"),
              dPR_mean=("delta_PR_AUC", "mean"),
              t_red_med=("t_reducer", "median"),
              t_mod_med=("t_model", "median"),
          )
          .reset_index()
          .sort_values(
              ["dataset", "model", "dAUC_mean"],
              ascending=[True, True, False],
          )
    )

    results = {
        "df": df,
        "agg": agg,
    }
    with open("results_B_real", "wb") as f:
        pickle.dump(results, f)

    return df, agg

\end{lstlisting}

\end{document}